% Generated mechanically from frozen input inputs/p09/ja/src/spaces-flat.tex.
% Do not edit this generated file; edit the bound locale source instead.

% OK, start here.
%

\setcounter{chapter}{76}
\chapter{代数空間上の平坦性}


\phantomsection
\label{spaces-flat-section-phantom}




\section{はじめに}
\label{spaces-flat-section-introduction}

\noindent
本章では、代数空間の枠組みにおける平坦加群および平坦射に関する
いくつかの高度な結果を論じる。まずスキームの枠組みにおける対応する章、
すなわち More on Flatness, Section \ref{flat-section-introduction}
に目を通すことを強く勧める。参考文献として、
Raynaud と Gruson による論文 \cite{GruRay} がある。







\section{不純性}
\label{spaces-flat-section-impure}

\noindent
本節は More on Flatness, Section \ref{flat-section-impure}
の類似である。

\begin{situation}
\label{spaces-flat-situation-pre-pure}
$S$ をスキームとする。$f : X \to Y$ を有限型かつ decent な射
\footnote{準分離射は decent である。次を参照せよ。
Decent Spaces, Lemma
\ref{decent-spaces-lemma-properties-trivial-implications}。
任意の射 $\Spec(k) \to Y$ で $k$ が体であるものに対して、代数空間
$X_k$ は $k$ 上有限表示である。実際、これは $k$ 上有限型であり、
Decent Spaces, Lemma
\ref{decent-spaces-lemma-locally-Noetherian-decent-quasi-separated}
により準分離だからである。}である $S$ 上の代数空間の射とする。
また、$\mathcal{F}$ を有限型の
準連接 $\mathcal{O}_X$-加群とする。最後に $y \in |Y|$ を $Y$ の点とする。
\end{situation}

\noindent
この状況で、スキーム $T$、射 $g : T \to Y$、
点 $t \in T$ で $g(t)=y$ を満たすもの、特殊化
$t' \leadsto t$（$T$ におけるもの）、および点
$\xi \in |X_T|$ で $t'$ の上にあるものを考える。
ここで $X_T = T \times_Y X$ である。図式は
\begin{equation}
\label{spaces-flat-equation-impurity}
\vcenter{
\xymatrix{
\xi \ar@{|->}[d] & \\
t' \ar@{~>}[r] & t \ar@{|->}[r] \ar[r] & y
}
}
\quad\quad
\vcenter{
\xymatrix{
X_T \ar[d]_{f_T} \ar[r] & X \ar[d]^f \\
T \ar[r]^g & Y
}
}
\end{equation}
である。さらに、$\mathcal{F}_T$ により
$\mathcal{F}$ の $X_T$ への引き戻しを表す。

\begin{definition}
\label{spaces-flat-definition-impurity}
Situation \ref{spaces-flat-situation-pre-pure} において、図式
(\ref{spaces-flat-equation-impurity}) が
{\it $\mathcal{F}$ の $y$ 上の不純性}を定めるとは、
$\xi \in \text{Ass}_{X_T/T}(\mathcal{F}_T)$ かつ
$t \not \in f_T(\overline{\{\xi\}})$ であることをいう。
これを「$(g : T \to Y, t' \leadsto t, \xi)$ を
$\mathcal{F}$ の $y$ 上の不純性とする」と表す。
\end{definition}

\noindent
同じことは次のようにも言える：
$(g : T \to Y, t' \leadsto t, \xi)$ が
$\mathcal{F}$ の $y$ 上の不純性であるとは、特殊化
$\xi \leadsto \theta$ であって位相空間 $|X_T|$ に属し、
$f_T(\theta)=t$ を満たすものが存在しないことをいう。
decent でない代数空間では、特殊化の挙動は良くない。射 $f$ が decent
ならば、$X_T$ は上のような任意の射 $g : T \to Y$ に対して
decent な代数空間である。次を参照せよ。
Decent Spaces, Definition \ref{decent-spaces-definition-relative-conditions}。

\begin{lemma}
\label{spaces-flat-lemma-impure-limit}
Situation \ref{spaces-flat-situation-pre-pure} において考える。
% P09-ERR-1236
$(g : T \to S, t' \leadsto t, \xi)$ を
$\mathcal{F}$ の $y$ 上の不純性とする。
$T = \lim_{i \in I} T_i$ が $Y$ 上のアフィンスキームの有向極限であると
仮定する。このとき、ある $i$ に対して三つ組
$(T_i \to Y, t'_i \leadsto t_i, \xi_i)$ は
$\mathcal{F}$ の $y$ 上の不純性である。
\end{lemma}

\begin{proof}
主張の記法は次の意味である。$p_i : T \to T_i$ を射影とし、
$t_i=p_i(t)$ および $t'_i=p_i(t')$ とおく。最後に
$\xi_i \in |X_{T_i}|$ を $\xi$ の像とする。
Divisors on Spaces, Lemma
\ref{spaces-divisors-lemma-base-change-relative-assassin}
により
$\xi_i \in \text{Ass}_{X_{T_i}/T_i}(\mathcal{F}_{T_i})$
である。したがって、示すべきことは
$t_i \not \in f_{T_i}(\overline{\{\xi_i\}})$ となる
$i$ が存在することだけである。

\medskip\noindent
$Z_i \subset X_{T_i}$ を
$\overline{\{\xi_i\}} \subset |X_{T_i}|$ 上の被約誘導スキーム構造とし、
$Z \subset X_T$ を
$\overline{\{\xi\}} \subset |X_T|$ 上の被約誘導スキーム構造とする。
すると
Limits of Spaces, Lemma \ref{spaces-limits-lemma-inverse-limit-irreducibles}
により $Z=\lim Z_i$ である
（各 $X_{T_i}$ が decent なので、この補題を適用できる）。
体 $k$ と射 $\Spec(k) \to T$ で像が $t$ となるものを選ぶ。このとき
$$
\emptyset =
Z \times_T \Spec(k) = (\lim Z_i) \times_{(\lim T_i)} \Spec(k)
= \lim Z_i \times_{T_i} \Spec(k)
$$
である。実際、極限はファイバー積と可換である
（極限どうしは可換である）。各
$Z_i \times_{T_i} \Spec(k)$ は準コンパクトである。なぜなら
$X_{T_i} \to T_i$ は有限型であり、したがって $Z_i \to T_i$ も
有限型だからである。よって
Limits of Spaces, Lemma \ref{spaces-limits-lemma-limit-nonempty}
により、$Z_i \times_{T_i} \Spec(k)$ が空となる $i$ が存在する。
合成 $\Spec(k) \to T \to T_i$ の像は $t_i$ なので、結論を得る。
\end{proof}

\noindent
不純性は平坦な基底変換に沿って上昇する。

\begin{lemma}
\label{spaces-flat-lemma-flat-ascent-impurity}
Situation \ref{spaces-flat-situation-pre-pure} において考える。
$(Y_1,y_1)\to(Y,y)$ を $S$ 上の点付き代数空間の射とする。
$Y_1\to Y$ は $y_1$ において平坦であると仮定する。
$(T \to Y, t' \leadsto t, \xi)$ が
$\mathcal{F}$ の $y$ 上の不純性ならば、不純性
$(T_1 \to Y_1, t_1' \leadsto t_1, \xi_1)$ であって、
$\mathcal{F}_1$（これは $\mathcal{F}$ の
$X_1=Y_1\times_YX$ への引き戻しである）の $y_1$ 上にあり、
$T_1$ が $Y_1\times_YT$ 上エタールとなるものが存在する。
\end{lemma}

\begin{proof}
エタール射 $T_1\to Y_1\times_YT$ で $T_1$ がスキームであるものを選び、
点 $t_1\in T_1$ で $y_1$ と $t$ に写るものをとる。
このような組 $(T_1,t_1)$ が存在することは
Properties of Spaces, Lemma \ref{spaces-properties-lemma-points-cartesian}
による。スキームの射 $T_1\to T$ は $t_1$ において平坦である
（Morphisms of Spaces, Lemma
\ref{spaces-morphisms-lemma-base-change-flat} と代数空間の平坦射の定義を用いる）。
% P09-ERR-1237
したがって特殊化 $t'_1\leadsto t_1$ で
$t'\leadsto t$ の上にあるものが存在する。次を参照せよ。
Morphisms, Lemma \ref{morphisms-lemma-generalizations-lift-flat}。
% P09-ERR-1238
点 $\xi_1\in|X_{T_1}|$ で $t'_1$ と $\xi$ に写り、かつ
$\xi_1 \in \text{Ass}_{X_{T_1}/T_1}(\mathcal{F}_{T_1})$
となるものを選ぶ。
$\Spec(\kappa(t'_1) \otimes_{\kappa(t')} \kappa(\xi))$ の点。
これは
Divisors on Spaces, Lemma
\ref{spaces-divisors-lemma-base-change-relative-assassin}
により可能である。$Z_1$ を $\{\xi_1\}$ の
$|X_{T_1}|$ における閉包とする。これは $\{\xi\}$ の
$|X_T|$ における閉包の中へ写るので、$Z_1$ の
$|T_1|$ における像は $t_1$ を含み得ない。
したがって
$(T_1 \to Y_1, t'_1 \leadsto t_1, \xi_1)$ は
% P09-ERR-1239
$\mathcal{F}_1$ の $Y_1$ 上の不純性である。
\end{proof}

\begin{lemma}
\label{spaces-flat-lemma-pure-along-X-y}
Situation \ref{spaces-flat-situation-pre-pure} において考える。
$\overline{y}$ を $y$ の上にある幾何学的点とする。
$\mathcal{O}=\mathcal{O}_{Y,\overline{y}}$ を
$Y$ の $\overline{y}$ におけるエタール局所環とする。
$Y^{sh}=\Spec(\mathcal{O})$、
$X^{sh}=X\times_YY^{sh}$ と書き、
$\mathcal{F}^{sh}$ を $\mathcal{F}$ の $X^{sh}$ への引き戻しとする。
次は同値である。
\begin{enumerate}
\item 不純性
$(Y^{sh} \to Y, y' \leadsto \overline{y}, \xi)$
が $\mathcal{F}$ の $y$ 上に存在する。
% P09-ERR-1240
\item $\text{Ass}_{X^{sh}/Y^{sh}}(\mathcal{F}^{sh})$ のすべての点は
閉ファイバー $X_{\overline{y}}$ の点へ特殊化する。
\item 不純性 $(T \to Y, t' \leadsto t, \xi)$ であって、
$\mathcal{F}$ の $y$ 上にあり、
$(T,t)\to(Y,y)$ がエタール近傍となるものが存在する。
\item 不純性 $(T \to Y, t' \leadsto t, \xi)$ であって、
$\mathcal{F}$ の $y$ 上にあり、
$T\to Y$ が $t$ において準有限となるものが存在する。
\end{enumerate}
\end{lemma}

\begin{proof}
(1) と (2) が同値であることは定義から直ちに従う。

\medskip\noindent
$\mathcal{O}=\mathcal{O}_{Y,\overline{y}}$ は、
$\mathcal{O}(V)$ の、エタール近傍
$(V,\overline{v})\to(Y,\overline{y})$ の圏上の
フィルター付き余極限であることを思い出そう
（Properties of Spaces, Lemma
\ref{spaces-properties-lemma-cofinal-etale}）。
さらに、アフィンなエタール近傍 $V$ だけを考えれば十分である。
したがって
$Y^{sh}=\Spec(\mathcal{O})=\lim\Spec(\mathcal{O}(V))=\lim V$
である。よって Lemma \ref{spaces-flat-lemma-impure-limit} により
(1) は (3) を含意する。

\medskip\noindent
エタール射は局所準有限であるから
（Morphisms of Spaces, Lemma
\ref{spaces-morphisms-lemma-etale-locally-quasi-finite}）、
(3) は (4) を含意する。

\medskip\noindent
最後に (4) を仮定する。$T$ を $t$ の開近傍で置き換えることにより、
$T\to Y$ が局所準有限であると仮定してよい。
Lemma \ref{spaces-flat-lemma-flat-ascent-impurity} により、不純性
$(T_1 \to Y^{sh}, t_1' \leadsto t_1, \xi_1)$ であって、
$T_1\to T\times_YY^{sh}$ がエタールとなるものを得る。
エタール射は局所準有限であり、さらに
Morphisms of Spaces, Lemma
\ref{spaces-morphisms-lemma-base-change-quasi-finite} および
Morphisms, Lemma \ref{morphisms-lemma-composition-quasi-finite}
を用いると、$T_1\to Y^{sh}$ は局所準有限である。
$\mathcal{O}$ は強ヘンゼル環なので、
More on Morphisms, Lemma
\ref{more-morphisms-lemma-etale-makes-quasi-finite-finite-at-point}
を適用できる。したがって $T_1$ を $t_1$ の開かつ閉な近傍で
置き換えることにより、
$T_1\to Y^{sh}=\Spec(\mathcal{O})$ が有限であると仮定してよい。
$\theta\in|X^{sh}|$ を $\xi_1$ の像とし、
$y'\in\Spec(\mathcal{O})$ を $t_1'$ の像とする。
Divisors on Spaces, Lemma
\ref{spaces-divisors-lemma-base-change-relative-assassin}
により
$\theta\in\text{Ass}_{X^{sh}/Y^{sh}}(\mathcal{F}^{sh})$ である。
$\pi:X_{T_1}\to X^{sh}$ は有限なので、
閉写像 $|X_{T_1}|\to|X^{sh}|$ を誘導する。したがって
$\overline{\{\xi_1\}}$ の像は $\overline{\{\theta\}}$ である。
よって $(Y^{sh}\to Y,y'\leadsto\overline{y},\theta)$ は
$\mathcal{F}$ の $y$ 上の不純性であり、証明は完了する。
\end{proof}







\section{相対的に純粋な加群}
\label{spaces-flat-section-pure}

\noindent
本節は More on Flatness, Section \ref{flat-section-pure}
の類似である。

\begin{definition}
\label{spaces-flat-definition-pure}
Situation \ref{spaces-flat-situation-pre-pure} において考える。
\begin{enumerate}
\item Lemma \ref{spaces-flat-lemma-pure-along-X-y} の同値な条件の
{\bf いずれも}成り立たないとき、$\mathcal{F}$ は
{\it $y$ 上で純粋}であるという。
\item $\mathcal{F}$ の $y$ 上の不純性が一つも存在しないとき、
$\mathcal{F}$ は {\it $y$ 上で普遍的に純粋}であるという。
\item $X$ が {\it $y$ 上で純粋}であるとは、
$\mathcal{O}_X$ が $y$ 上で純粋であることをいう。
\item $\mathcal{F}$ が {\it 普遍的に $Y$-純粋}、または
{\it $Y$ に関して普遍的に純粋}であるとは、$\mathcal{F}$ が
$y$ 上で普遍的に純粋であることがすべての $y\in|Y|$ に対して
成り立つことをいう。
\item $\mathcal{F}$ が {\it $Y$-純粋}、または
{\it $Y$ に関して純粋}であるとは、$\mathcal{F}$ が
$y$ 上で純粋であることがすべての $y\in|Y|$ に対して
成り立つことをいう。
\item $X$ が {\it $Y$-純粋}、または
{\it $Y$ に関して純粋}であるとは、$\mathcal{O}_X$ が
$Y$ に関して純粋であることをいう。
\end{enumerate}
\end{definition}

\noindent
必要となる補題を述べる。

\begin{lemma}
\label{spaces-flat-lemma-base-change-universally}
Situation \ref{spaces-flat-situation-pre-pure} において考える。
% P09-ERR-1241
\begin{enumerate}
\item $\mathcal{F}$ は $y$ 上で普遍的に純粋である。
\item 点付き代数空間の任意の射 $(Y',y')\to(Y,y)$ に対して、
引き戻し $\mathcal{F}_{Y'}$ は $y'$ 上で純粋である。
\end{enumerate}
特に、$\mathcal{F}$ が $Y$ に関して普遍的に純粋であるための
必要十分条件は、すべての基底変換 $\mathcal{F}_{Y'}$、すなわち
$\mathcal{F}$ の基底変換が $Y'$ に関して純粋であることである。
\end{lemma}

\begin{proof}
これは形式的である。
\end{proof}

\begin{lemma}
\label{spaces-flat-lemma-quasi-finite-base-change}
Situation \ref{spaces-flat-situation-pre-pure} において考える。
$(Y',y')\to(Y,y)$ を点付き代数空間の射とする。
$Y'\to Y$ が $y'$ において準有限であり、$\mathcal{F}$ が
$y$ 上で純粋ならば、$\mathcal{F}_{Y'}$ は $y'$ 上で純粋である。
\end{lemma}

\begin{proof}
% P09-ERR-1242
$(T\to Y',t'\leadsto t,\xi)$ が $\mathcal{F}_{Y'}$ の
$y'$ 上の不純性であり、$T\to Y'$ が $t$ において
準有限ならば、$(T\to Y,t'\to t,\xi)$ は $\mathcal{F}$ の
$y$ 上の不純性であり、$T\to Y$ は $t$ において準有限である。
次を参照せよ。
Morphisms of Spaces, Lemma
\ref{spaces-morphisms-lemma-composition-quasi-finite}。
したがって、この補題は純粋性の定義から直ちに従う。
\end{proof}

\noindent
純粋性は平坦降下を満たす。

\begin{lemma}
\label{spaces-flat-lemma-flat-descend-pure}
Situation \ref{spaces-flat-situation-pre-pure} において考える。
$(Y_1,y_1)\to(Y,y)$ を点付き代数空間の射とする。
$Y_1\to Y$ は $y_1$ において平坦であると仮定する。
\begin{enumerate}
\item $\mathcal{F}_{Y_1}$ が $y_1$ 上で純粋ならば、
$\mathcal{F}$ は $y$ 上で純粋である。
\item $\mathcal{F}_{Y_1}$ が $y_1$ 上で普遍的に純粋ならば、
$\mathcal{F}$ は $y$ 上で普遍的に純粋である。
\end{enumerate}
\end{lemma}

\begin{proof}
これは、不純性が平坦な基底変換に沿って上昇することによる。
Lemma \ref{spaces-flat-lemma-flat-ascent-impurity} を参照せよ。
例えば (1) は次のように従う。
% P09-ERR-1243
任意の不純性 $(T\to Y,t'\leadsto t,\xi)$ であって
$\mathcal{F}$ の $y$ 上にあり、$T\to Y$ が $t$ において
準有限であるものから、この補題により、不純性
$(T_1\to Y_1,t_1'\leadsto t_1,\xi_1)$ が得られる。
これは、引き戻し $\mathcal{F}_1$、すなわち $\mathcal{F}$ の
$X_1=Y_1\times_YX$ への引き戻しの $y_1$ 上にあり、
$T_1$ は $Y_1\times_YT$ 上エタールである。
したがって $T_1\to Y_1$ は $t_1$ において準有限である。
実際、エタール射は局所準有限であり、局所準有限射の合成は
局所準有限である
（Morphisms of Spaces, Lemmas
\ref{spaces-morphisms-lemma-etale-locally-quasi-finite} および
\ref{spaces-morphisms-lemma-composition-quasi-finite}）。
(2) についても同様である。
\end{proof}

\begin{lemma}
\label{spaces-flat-lemma-supported-on-closed}
Situation \ref{spaces-flat-situation-pre-pure} において考える。
$i:Z\to X$ を閉埋め込みとし、
$\mathcal{F}=i_*\mathcal{G}$ であると仮定する。ここで
$\mathcal{G}$ は $Z$ 上の有限型準連接層である。
このとき、$\mathcal{G}$ が $y$ 上で（普遍的に）純粋であるための
必要十分条件は、$\mathcal{F}$ が $y$ 上で（普遍的に）
純粋であることである。
\end{lemma}

\begin{proof}
これは Divisors on Spaces, Lemma
\ref{spaces-divisors-lemma-relative-weak-assassin-finite}
から従う。
\end{proof}

\begin{lemma}
\label{spaces-flat-lemma-proper-pure}
Situation \ref{spaces-flat-situation-pre-pure} において考える。
\begin{enumerate}
\item $\mathcal{F}$ の台が $Y$ 上固有ならば、
$\mathcal{F}$ は $Y$ に関して普遍的に純粋である。
\item $f$ が固有ならば、$\mathcal{F}$ は
$Y$ に関して普遍的に純粋である。
\item $f$ が固有ならば、$X$ は $Y$ に関して普遍的に純粋である。
\end{enumerate}
\end{lemma}

\begin{proof}
まず (1) を (2) に帰着する。すなわち、
$Z\subset X$ を $\mathcal{F}$ のスキーム論的台とする
（Morphisms of Spaces, Definition
\ref{spaces-morphisms-definition-scheme-theoretic-support}）。
$i:Z\to X$ を対応する閉埋め込みとし、
$\mathcal{F}=i_*\mathcal{G}$ と書く。ここで有限型準連接
$\mathcal{O}_Z$-加群 $\mathcal{G}$ を用いた。
(1) の場合、仮定により
$Z\to Y$ は固有である。したがって
Lemma \ref{spaces-flat-lemma-supported-on-closed} により、(1) は (2) に帰着される。

\medskip\noindent
$f$ が固有であると仮定する。
$(g:T\to Y,t'\leadsto t,\xi)$ を
$\mathcal{F}$ の $y$ 上の不純性とする。
$f$ は固有なので普遍閉であり、したがって
$f_T:X_T\to T$ は閉写像である。$f_T(\xi)=t'$ だから、
% P09-ERR-1244
これは $t\in f(\overline{\{\xi\}})$ を含意し、矛盾である。
\end{proof}







\section{有限型平坦加群}
\label{spaces-flat-section-finite-type-flat}

\noindent
More on Flatness, Sections
\ref{flat-section-finite-type-flat-I},
\ref{flat-section-finite-type-flat-II}, and
\ref{flat-section-finite-type-flat-III}
と比較されたい。
% P09-ERR-1245
これらの結果の大部分は、エタール局所化により代数空間の
直接の帰結をもつ。

\begin{lemma}
\label{spaces-flat-lemma-existence-complete}
$S$ をスキームとする。
$X\to Y$ を $S$ 上の代数空間の有限型射とする。
$\mathcal{F}$ を有限型準連接 $\mathcal{O}_X$-加群とする。
$y\in|Y|$ を点とする。このとき、エタール射
$(Y',y')\to(Y,y)$ で $Y'$ がアフィンスキームであるものと、エタール射
$h_i:W_i\to X_{Y'}$（$i=1,\ldots,n$）で、各 $i$ に対して
$\mathcal{F}_i/W_i/Y'$ の $y'$ 上の完全なデヴィサージュが存在するものが
存在する。ここで $\mathcal{F}_i$ は $\mathcal{F}$ の $W_i$ への
引き戻しであり、さらに
$|(X_{Y'})_{y'}|\subset\bigcup h_i(W_i)$ が成り立つ。
\end{lemma}

\begin{proof}
問題は $Y$ 上エタール局所的なので、$Y$ がアフィンスキームであると
仮定してよい。このとき $X$ は準コンパクトであるから、
アフィンスキーム $X'$ と全射エタール射 $X'\to X$ を選べる。
そこで $X'\to Y$、$(X'\to Y)^*\mathcal{F}$、および $y$ に
More on Flatness, Lemma \ref{flat-lemma-existence-complete}
を適用すれば、所望の結果を得る。
\end{proof}

\begin{lemma}
\label{spaces-flat-lemma-open-in-fibre-where-flat}
$S$ をスキームとする。
$f:X\to Y$ を $S$ 上の代数空間の局所有限型射とする。
$\mathcal{F}$ を有限型準連接 $\mathcal{O}_X$-加群とする。
$y\in|Y|$ とし、$F=f^{-1}(\{y\})\subset|X|$ とおく。このとき集合
$$
\{x \in F \mid \mathcal{F} \text{ は }Y\text{ 上、次の点で平坦：}x\}
$$
は $F$ において開である。
\end{lemma}

\begin{proof}
スキーム $V$、点 $v\in V$、およびエタール射
$V\to Y$ で $v$ を $y$ に写すものを選ぶ。スキーム $U$ と全射エタール射
$U\to V\times_YX$ を選ぶ。$|U_v|\to F$ は位相空間の開連続写像である。
実際、$|U|\to|X|$ は連続かつ開である。したがって結果は
スキームの場合、すなわち More on Flatness, Lemma
\ref{flat-lemma-open-in-fibre-where-flat} から従う。
\end{proof}

\begin{lemma}
\label{spaces-flat-lemma-bourbaki-finite-type-general-base-at-point}
$S$ をスキームとする。
$f:X\to Y$ を $S$ 上の代数空間の局所有限型射とする。
$x\in|X|$ とし、その像を $y\in|Y|$ とする。
$\mathcal{F}$ を $X$ 上の有限型準連接層とし、
$\mathcal{G}$ を $Y$ 上の準連接層とする。
$\mathcal{F}$ が $x$ において $Y$ 上平坦ならば、
$$
x \in \text{WeakAss}_X(\mathcal{F} \otimes_{\mathcal{O}_X} f^*\mathcal{G})
\Leftrightarrow
y \in \text{WeakAss}_Y(\mathcal{G})
\text{ and }
x \in \text{Ass}_{X/Y}(\mathcal{F}).
$$
\end{lemma}

\begin{proof}
可換図式
$$
\xymatrix{
U \ar[d] \ar[r]_g & V \ar[d] \\
X \ar[r]^f & Y
}
$$
で、$U$ と $V$ がスキームであり、縦の矢印が全射エタールであるものを選ぶ。
$u\in U$ を選び、これが $x$ に写るようにする。
$\mathcal{E}=\mathcal{F}|_U$ および
$\mathcal{H}=\mathcal{G}|_V$ とおく。$v\in V$ を $u$ の像とする。
このとき
$x\in\text{WeakAss}_X(\mathcal{F}\otimes_{\mathcal{O}_X}f^*\mathcal{G})$
であることと
% P09-ERR-1246
$u\in\text{WeakAss}_X(\mathcal{E}\otimes_{\mathcal{O}_X}g^*\mathcal{H})$
であることは、
Divisors on Spaces, Definition
\ref{spaces-divisors-definition-weakly-associated}
により同値である。同様に、
$y\in\text{WeakAss}_Y(\mathcal{G})$ であることと
$v\in\text{WeakAss}_V(\mathcal{H})$ であることは同値である。
最後に、Divisors on Spaces, Definition
\ref{spaces-divisors-definition-relative-weak-assassin}
により、$x\in\text{Ass}_{X/Y}(\mathcal{F})$ であることと
$u\in\text{Ass}_{U_v}(\mathcal{E}|_{U_v})$ であることは同値である。
$\mathcal{F}$ の $x$ における平坦性は
$\mathcal{E}$ の $u$ における平坦性と同値であることに注意せよ。
Morphisms of Spaces, Definition
\ref{spaces-morphisms-definition-flat-module} を参照せよ。
$g:U\to V$、$\mathcal{E}$、$\mathcal{H}$、$u$、$v$ に対する同値は
More on Flatness, Lemma
\ref{flat-lemma-bourbaki-finite-type-general-base-at-point}
である。
\end{proof}

\begin{lemma}
\label{spaces-flat-lemma-bourbaki-finite-type-general-base}
$S$ をスキームとする。$f:X\to Y$ を $S$ 上の代数空間の
局所有限型射とする。$\mathcal{F}$ を $X$ 上の有限型準連接層で
$Y$ 上平坦なものとし、$\mathcal{G}$ を $Y$ 上の準連接層とする。
このとき
$$
\text{WeakAss}_X(\mathcal{F} \otimes_{\mathcal{O}_X} f^*\mathcal{G}) =
\text{Ass}_{X/Y}(\mathcal{F}) \cap
|f|^{-1}(\text{WeakAss}_Y(\mathcal{G}))
$$
である。
\end{lemma}

\begin{proof}
Lemma \ref{spaces-flat-lemma-bourbaki-finite-type-general-base-at-point}
の直ちの帰結である。
\end{proof}

\begin{theorem}
\label{spaces-flat-theorem-finite-type-flat}
$S$ をスキームとする。
$f:X\to Y$ を $S$ 上の代数空間の射とし、
$\mathcal{F}$ を準連接 $\mathcal{O}_X$-加群とする。次を仮定する。
\begin{enumerate}
\item $X\to Y$ は局所有限表示である。
\item $\mathcal{F}$ は有限型 $\mathcal{O}_X$-加群である。
\item $Y$ の弱随伴点の集合は $Y$ において局所有限である。
\end{enumerate}
このとき
$U=\{x\in|X|:\mathcal{F}\text{ は次の点で平坦：}x\text{、基底：}Y\}$
は $X$ において開であり、$\mathcal{F}|_U$ は有限表示
$\mathcal{O}_U$-加群であって $Y$ 上平坦である。
\end{theorem}

\begin{proof}
条件 (3) は、全射エタール射 $V\to Y$ で $V$ がスキームであるものに対して、
$V$ の弱随伴点がスキーム $V$ 上で局所有限であることを意味する
（Divisors on Spaces, Definition
\ref{spaces-divisors-definition-weakly-associated}
により、$V$ の弱随伴点はちょうど $Y$ の弱随伴点の逆像であることを
思い出そう）。以上を踏まえると、問題は $X$ と $Y$ 上エタール局所的で
あるから、$X$ と $Y$ がスキームであると仮定してよい。
したがって結果は More on Flatness, Theorem
\ref{flat-theorem-finite-type-flat} から従う。
\end{proof}

\begin{lemma}
\label{spaces-flat-lemma-finite-type-flat-along-fibre-free-variant}
$S$ をスキームとする。$f:X\to Y$ を $S$ 上の代数空間の射とし、
$\mathcal{F}$ を $X$ 上の準連接層とする。
$y\in|Y|$ とし、$F=f^{-1}(\{y\})\subset|X|$ とおく。次を仮定する。
\begin{enumerate}
\item $f$ は有限型である。
\item $\mathcal{F}$ は有限型である。
\item $\mathcal{F}$ は $Y$ 上で、すべての $x\in F$ において平坦である。
\end{enumerate}
このとき、エタール射 $(Y',y')\to(Y,y)$ で
$Y'$ がスキームであるものと、代数空間の可換図式
$$
\xymatrix{
X \ar[d] & X' \ar[l]^g \ar[d] \\
Y & \Spec(\mathcal{O}_{Y', y'}) \ar[l]
}
$$
であって、$X'\to X\times_Y\Spec(\mathcal{O}_{Y',y'})$ がエタール、
$|X'_{y'}|\to F$ が全射、$X'$ がアフィン、かつ
$\Gamma(X',g^*\mathcal{F})$ が自由
$\mathcal{O}_{Y',y'}$-加群となるものが存在する。
\end{lemma}

\begin{proof}
エタール射 $(Y',y')\to(Y,y)$ で
$Y'$ がアフィンスキームであるものを選ぶ。このとき
$X\times_YY'$ は準コンパクトである。
アフィンスキーム $X'$ と全射エタール射
$X'\to X\times_YY'$ を選ぶ。図式は
$$
\xymatrix{
X \ar[d] & X' \ar[l]^g \ar[d] \\
Y & Y' \ar[l]
}
$$
である。すると $\mathcal{F}'=g^*\mathcal{F}$ は
$Y'$ 上で、$X'_{y'}$ のすべての点において平坦である。
Morphisms of Spaces, Lemma
\ref{spaces-morphisms-lemma-base-change-module-flat}
を参照せよ。したがって、射 $X'\to Y'$、準連接層
$g^*\mathcal{F}$、および点 $y'$ にスキームの場合の補題
（More on Flatness, Lemma
\ref{flat-lemma-finite-type-flat-along-fibre-free-variant}）
を適用できる。これにより、エタール射
$(Y'',y'')\to(Y',y')$ と可換図式
$$
\xymatrix{
X \ar[d] & X' \ar[l]^g \ar[d] & X'' \ar[l]^{g'} \ar[d] \\
Y & Y' \ar[l] & \Spec(\mathcal{O}_{Y'', y''}) \ar[l]
}
$$
を得る。所望のものとして $(Y'',y'')\to(Y,y)$ と
$g\circ g':X''\to X$ をとればよい。
\end{proof}

\begin{theorem}
\label{spaces-flat-theorem-check-flatness-at-associated-points}
$S$ をスキームとする。
$f:X\to Y$ を $S$ 上の代数空間の局所有限型射とし、
$\mathcal{F}$ を有限型準連接 $\mathcal{O}_X$-加群とする。
$x\in|X|$ とし、その像を $y\in|Y|$ とする。
$F=f^{-1}(\{y\})\subset|X|$ とおく。次の条件を考える。
\begin{enumerate}
\item $\mathcal{F}$ は $x$ において $Y$ 上平坦である。
\item 任意の $x'\in F\cap\text{Ass}_{X/Y}(\mathcal{F})$ で
$x$ へ特殊化するものに対して、
$\mathcal{F}$ は $x'$ において $Y$ 上平坦である。
\end{enumerate}
常に (2) $\Rightarrow$ (1) が成り立つ。
$X$ と $Y$ が decent ならば、(1) $\Rightarrow$ (2) も成り立つ。
\end{theorem}

\begin{proof}
(2) を仮定する。スキーム $V$ と全射エタール射 $V\to Y$ を選ぶ。
スキーム $U$ と全射エタール射 $U\to V\times_YX$ を選ぶ。
点 $u\in U$ で $x$ に写るものを選び、$v\in V$ を $u$ の像とする。
% P09-ERR-1247
$\mathcal{F}|_U=(U\to X)^*\mathcal{F}$ と点 $u$ に対する対応する結果から
結論を導く。$U_v$。
これは、$\text{Ass}_{U/V}(\mathcal{F}|_U)\cap|U_v|$ が
$\text{Ass}_{U_v}(\mathcal{F}|_{U_v})$ に等しく、さらに
$F\cap\text{Ass}_{X/Y}(\mathcal{F})$ の逆像に等しいためである。
写像 $|U_v|\to F$ は連続なので、$|U_v|$ における特殊化は
$F$ における特殊化へ写る。したがって条件 (2) は
$U\to V$、$\mathcal{F}|_U$、および点 $u$ に継承される。
よって More on Flatness, Theorem
\ref{flat-theorem-check-flatness-at-associated-points}
を適用でき、(1) が成り立つ。

\medskip\noindent
$Y$ が decent ならば、$y$ を準コンパクトな単射
$\Spec(k)\to Y$ によって表せる
（decent 空間の定義による。Decent Spaces, Definition
\ref{decent-spaces-definition-very-reasonable} を参照せよ）。
このとき
Decent Spaces, Lemma \ref{decent-spaces-lemma-topology-fibre}
により $F=|X_k|$ である。さらに $X$ も decent ならば
（より一般には $f$ が decent ならばよい。Decent Spaces, Definition
\ref{decent-spaces-definition-relative-conditions} および
Decent Spaces, Lemma
\ref{decent-spaces-lemma-property-for-morphism-out-of-property}
を参照せよ）、$X_y$ も decent 空間である。
さらに、$F$ における特殊化は $U_v\to X_y$ における
特殊化へ持ち上げられる。次を参照せよ。
Decent Spaces, Lemma
\ref{decent-spaces-lemma-decent-specialization}。
以上を踏まえれば、逆の含意も成り立つことは明らかである。
実際、これはスキームの場合に成り立つ。
\end{proof}

\begin{lemma}
\label{spaces-flat-lemma-check-along-closed-fibre}
$S$ を閉点 $s$ をもつ局所スキームとする。
$f:X\to S$ を代数空間 $X$ から $S$ への局所有限型射とする。
$\mathcal{F}$ を有限型準連接 $\mathcal{O}_X$-加群とする。次を仮定する。
\begin{enumerate}
\item $\text{Ass}_{X/S}(\mathcal{F})$ のすべての点は
閉ファイバー $X_s$ の点へ特殊化する\footnote{例えば、
$f$ が有限型で $\mathcal{F}$ が $X_s$ に沿って純粋である場合、
または $f$ が固有である場合にこれは成り立つ。}。
\item $\mathcal{F}$ は $S$ 上で、$X_s$ のすべての点において平坦である。
\end{enumerate}
このとき $\mathcal{F}$ は $S$ 上平坦である。
\end{lemma}

\begin{proof}
Theorem \ref{spaces-flat-theorem-check-flatness-at-associated-points}
により、$\mathcal{F}$ の $S$ 上の相対随伴点において
平坦性を確認すれば十分であることから直ちに従う。
\end{proof}




\section{有限表示平坦加群}
\label{spaces-flat-section-finitely-presented-flat}

\noindent
これは More on Flatness, Section
\ref{flat-section-finitely-presented-flat}
の類似である。

\begin{proposition}
\label{spaces-flat-proposition-finite-presentation-flat-at-point}
$S$ をスキームとする。
$f:X\to Y$ を $S$ 上の代数空間の射とする。
$\mathcal{F}$ を $X$ 上の準連接層とする。
$x\in|X|$ とし、その像を $y\in|Y|$ とする。次を仮定する。
\begin{enumerate}
\item $f$ は局所有限表示である。
\item $\mathcal{F}$ は有限表示である。
\item $\mathcal{F}$ は $x$ において $Y$ 上平坦である。
\end{enumerate}
このとき、点付きスキームの可換図式
$$
\xymatrix{
(X, x) \ar[d] & (X', x') \ar[l]^g \ar[d] \\
(Y, y) & (Y', y') \ar[l]
}
$$
であって、横の矢印がエタールであり、$X'$ と $Y'$ がアフィンで、
$\Gamma(X',g^*\mathcal{F})$ が射影
$\Gamma(Y',\mathcal{O}_{Y'})$-加群となるものが存在する。
\end{proposition}

\begin{proof}
この形で述べれば、この命題は直ちにスキームの場合、すなわち
More on Flatness, Proposition
\ref{flat-proposition-finite-presentation-flat-at-point}
に帰着される。
\end{proof}

\begin{lemma}
\label{spaces-flat-lemma-finite-presentation-flat-along-fibre}
$S$ をスキームとする。
$f:X\to Y$ を $S$ 上の代数空間の射とする。
$\mathcal{F}$ を $X$ 上の準連接層とする。
$y\in|Y|$ とし、$F=f^{-1}(\{y\})\subset|X|$ とおく。次を仮定する。
\begin{enumerate}
\item $f$ は有限表示である。
\item $\mathcal{F}$ は有限表示である。
\item $\mathcal{F}$ は $Y$ 上で、すべての $x\in F$ において平坦である。
\end{enumerate}
このとき、代数空間の可換図式
$$
\xymatrix{
X \ar[d] & X' \ar[l]^g \ar[d] \\
Y & Y' \ar[l]_h
}
$$
であって、$h$ と $g$ がエタールであり、点
$y'\in|Y'|$ で $y$ に写るものが存在し、$F\subset g(|X'|)$ が成り立ち、
代数空間 $X'$ と $Y'$ がアフィンで、かつ
$\Gamma(X',g^*\mathcal{F})$ が射影
$\Gamma(Y',\mathcal{O}_{Y'})$-加群となるものが存在する。
\end{lemma}

\begin{proof}
この形で述べれば、この補題は直ちにスキームの場合、すなわち
More on Flatness, Lemma
\ref{flat-lemma-finite-presentation-flat-along-fibre}
に帰着される。
\end{proof}




\section{純粋性の判定条件}
\label{spaces-flat-section-criterion-purity}

\noindent
本節は More on Flatness, Section
\ref{flat-section-criterion-purity}
の類似である。

\begin{lemma}
\label{spaces-flat-lemma-associated-point-specializes}
$S$ をスキームとする。$X$ を $S$ 上局所有限型な decent 代数空間とする。
$\mathcal{F}$ を有限型準連接 $\mathcal{O}_X$-加群とする。
$s\in S$ を、$\mathcal{F}$ が $S$ 上で $X_s$ のすべての点において
平坦となるような点とする。
$x'\in\text{Ass}_{X/S}(\mathcal{F})$ とする。
$\{x'\}$ の $|X|$ における閉包が $|X_s|$ と交わるならば、その閉包は
$\text{Ass}_{X/S}(\mathcal{F})\cap|X_s|$ と交わる。
\end{lemma}

\begin{proof}
$|X_s|\subset|X|$ は、$|X|$ の点のうち $s\in S$ の上にあるものの集合で
あることに注意せよ。次を参照せよ。
Decent Spaces, Lemma \ref{decent-spaces-lemma-topology-fibre}。
$t\in|X_s|$ を、$x'$ の $|X|$ における特殊化とする。
アフィンスキーム $U$、点 $u\in U$、およびエタール射
$\varphi:U\to X$ で $u$ を $t$ に写すものを選ぶ。
Decent Spaces, Lemma
\ref{decent-spaces-lemma-decent-specialization}
により、特殊化 $u'\leadsto u$ で $u'$ が $x'$ に写るものを選べる。
$g=f\circ\varphi$ とおく。
$s'=g(u')=f(x')$ は $s$ へ特殊化することに注意せよ。
$\text{Ass}_{X/S}(\mathcal{F})$ の定義により
$u'\in\text{Ass}_{U/S}(\varphi^*\mathcal{F})$ である。
この補題のスキーム版
（More on Flatness, Lemma
\ref{flat-lemma-associated-point-specializes}）
により、特殊化 $u'\leadsto u$ で
$$
u \in \text{Ass}_{U_s}(\varphi^*\mathcal{F}_s) =
\text{Ass}_{U/S}(\varphi^*\mathcal{F}) \cap U_s
$$
となるものが存在する。したがって
$x=\varphi(u)\in\text{Ass}_{X/S}(\mathcal{F})$ は $s$ の上にあり、
補題が証明された。
\end{proof}

\begin{lemma}
\label{spaces-flat-lemma-contains-relative-ass-after-base-change}
$Y$ をスキーム $S$ 上の代数空間とする。
$g:X'\to X$ を $Y$ 上の代数空間の射とし、$X$ は $Y$ 上局所有限型であると
する。$\mathcal{F}$ を準連接 $\mathcal{O}_X$-加群とする。
$\text{Ass}_{X/Y}(\mathcal{F})\subset g(|X'|)$ ならば、任意の射
$Z\to Y$ に対して
$\text{Ass}_{X_Z/Z}(\mathcal{F}_Z)\subset g_Z(|X'_Z|)$
が成り立つ。
\end{lemma}

\begin{proof}
Properties of Spaces, Lemma
\ref{spaces-properties-lemma-points-cartesian}
により、写像
$|X'_Z|\to|X_Z|\times_{|X|}|X'|$ は全射である。実際、
$X'_Z$ は $X_Z\times_XX'$ に等しい。
Divisors on Spaces, Lemma
\ref{spaces-divisors-lemma-base-change-relative-assassin}
により、写像 $|X_Z|\to|X|$ は
$\text{Ass}_{X_Z/Z}(\mathcal{F}_Z)$ を
$\text{Ass}_{X/Y}(\mathcal{F})$ の中へ写す。これで補題が従う。
\end{proof}

\begin{lemma}
\label{spaces-flat-lemma-pure-on-top}
$Y$ をスキーム $S$ 上の代数空間とする。
$g:X'\to X$ を $Y$ 上の代数空間のエタール射とする。
構造射 $X'\to Y$ および $X\to Y$ は decent かつ有限型であると仮定する。
$\mathcal{F}$ を有限型準連接 $\mathcal{O}_X$-加群とする。
$y\in|Y|$ とし、% P09-ERR-1248
$F=f^{-1}(\{y\})\subset|X|$ とおく。
\begin{enumerate}
\item $\text{Ass}_{X/Y}(\mathcal{F})\subset g(|X'|)$ であり、
$g^*\mathcal{F}$ が $y$ 上で（普遍的に）純粋ならば、
$\mathcal{F}$ は $y$ 上で（普遍的に）純粋である。
\item $\mathcal{F}$ が $y$ 上で純粋であり、$g(|X'|)$ が $F$ を含み、
$Y$ が閉点 $y$ をもつアフィン局所空間ならば、
$\text{Ass}_{X/Y}(\mathcal{F})\subset g(|X'|)$ である。
\item $\mathcal{F}$ が $y$ 上で純粋であり、$\mathcal{F}$ が
$F$ のすべての点で平坦で、$g(|X'|)$ が
$\text{Ass}_{X/Y}(\mathcal{F})\cap F$ を含み、$Y$ が閉点 $y$ をもつ
アフィン局所空間ならば、
$\text{Ass}_{X/Y}(\mathcal{F})\subset g(|X'|)$ である。
\item % P09-ERR-1249
ここにさらに追加する。
\end{enumerate}
\end{lemma}

\begin{proof}
$X\to Y$ と $X'\to Y$ に関する仮定により、Sections
\ref{spaces-flat-section-impure} および \ref{spaces-flat-section-pure}
の内容をこれらの射と層 $\mathcal{F}$ および
$g^*\mathcal{F}$ に適用できる。$g$ はエタールなので、
$\text{Ass}_{X'/Y}(g^*\mathcal{F})$ は
$\text{Ass}_{X/Y}(\mathcal{F})$ の逆像であり、基底変換後にも同じことが
成り立つ。

\medskip\noindent
(1) の証明。
$\text{Ass}_{X/Y}(\mathcal{F})\subset g(|X'|)$ と仮定する。
$(T\to Y,t'\leadsto t,\xi)$ が $\mathcal{F}$ の $y$ 上の不純性であるとする。
Lemma \ref{spaces-flat-lemma-contains-relative-ass-after-base-change} により
$\text{Ass}_{X_T/T}(\mathcal{F}_T)\subset g_T(|X'_T|)$
であるから、点 $\xi'\in|X'_T|$ で $\xi$ に写るものを選べる。
上で述べたことにより、
$(T\to Y,t'\leadsto t,\xi')$ は $g^*\mathcal{F}$ の
% P09-ERR-1250
$y'$ 上の不純性である。これより (1) が成り立つ。

\medskip\noindent
(2) の証明。$g(|X'|)$ が $|X|$ において開であること、および純粋性により
$\text{Ass}_{X/Y}(\mathcal{F})$ のすべての点が $F$ の点へ特殊化すること
から従う。

\medskip\noindent
(3) の証明。$g(|X'|)$ が $|X|$ において開であること、および純粋性と
Lemma \ref{spaces-flat-lemma-associated-point-specializes} により
$\text{Ass}_{X/Y}(\mathcal{F})$ のすべての点が
$\text{Ass}_{X/Y}(\mathcal{F})\cap F$ の点へ特殊化することから従う。
\end{proof}

\begin{lemma}
\label{spaces-flat-lemma-finite-type-flat-pure-along-fibre-is-universal}
$S$ をスキームとする。$f:X\to Y$ を $S$ 上の代数空間の射とする。
$\mathcal{F}$ を準連接 $\mathcal{O}_X$-加群とする。
$y\in|Y|$ とする。次を仮定する。
\begin{enumerate}
\item $f$ は decent かつ有限型である。
\item $\mathcal{F}$ は有限型である。
\item $\mathcal{F}$ は $Y$ 上で、$y$ の上にあるすべての点において
平坦である。
\item $\mathcal{F}$ は $y$ 上で純粋である。
\end{enumerate}
このとき $\mathcal{F}$ は $y$ 上で普遍的に純粋である。
\end{lemma}

\begin{proof}
射 $\Spec(\mathcal{O}_{Y,\overline{y}})\to Y$ を考える。
これは強ヘンゼル局所環のスペクトルからの平坦射であり、閉点を
$y$ に写す。Lemma \ref{spaces-flat-lemma-flat-descend-pure} により、
次の段落で記述する場合に帰着される。

\medskip\noindent
$Y$ が強ヘンゼル局所環 $R$ のスペクトルであり、
閉点が $y$ であると仮定する。
Lemma \ref{spaces-flat-lemma-finite-type-flat-along-fibre-free-variant}
により、エタール射 $g:X'\to X$ で
$g(|X'|)\supset|X_y|$、$X'$ がアフィン、かつ
$\Gamma(X',g^*\mathcal{F})$ が自由 $R$-加群となるものが存在する。
すると $g^*\mathcal{F}$ は $Y$ に関して普遍的に純粋である。
次を参照せよ。
More on Flatness, Lemma
\ref{flat-lemma-affine-locally-projective-pure}。
したがって
$g(|X'|)$ が $\text{Ass}_{X/Y}(\mathcal{F})$ を含むことを示せば十分である。
Lemma \ref{spaces-flat-lemma-pure-on-top} の (1) を参照せよ。
これはさらに Lemma \ref{spaces-flat-lemma-pure-on-top} の (2) から従う。
\end{proof}

\begin{lemma}
\label{spaces-flat-lemma-finite-type-flat-pure-is-universal}
$S$ をスキームとする。
$f:X\to Y$ を $S$ 上の代数空間の decent な有限型射とする。
$\mathcal{F}$ を有限型準連接 $\mathcal{O}_X$-加群とする。
$\mathcal{F}$ は $Y$ 上平坦であると仮定する。この場合、
$\mathcal{F}$ が $Y$ に関して純粋であるための必要十分条件は、
$\mathcal{F}$ が $Y$ に関して普遍的に純粋であることである。
\end{lemma}

\begin{proof}
Lemma \ref{spaces-flat-lemma-finite-type-flat-pure-along-fibre-is-universal}
と定義の直ちの帰結である。
\end{proof}

\begin{lemma}
\label{spaces-flat-lemma-universally-separating}
$Y$ をスキーム $S$ 上の代数空間とする。
$g:X'\to X$ を $Y$ 上の代数空間の平坦射とし、
$X$ は $Y$ 上局所有限型であるとする。
$\mathcal{F}$ を有限型準連接 $\mathcal{O}_X$-加群で
$Y$ 上平坦なものとする。
$\text{Ass}_{X/Y}(\mathcal{F})\subset g(|X'|)$ ならば、標準写像
$$
\mathcal{F} \longrightarrow g_*g^*\mathcal{F}
$$
は単射であり、任意の基底変換後にも単射のままである。
\end{lemma}

\begin{proof}
最後の主張は、$\mathcal{F}_Z\to(g_Z)_*g_Z^*\mathcal{F}_Z$ が
任意の射 $Z\to Y$ に対して単射であるという意味である。
相対随伴点に関する仮定は基底変換で保たれるから
（Lemma \ref{spaces-flat-lemma-contains-relative-ass-after-base-change}）、
表示された矢印の単射性を示せば十分である。

\medskip\noindent
$\mathcal{K}=\Ker(\mathcal{F}\to g_*g^*\mathcal{F})$ とおく。
目標は $\mathcal{K}=0$ を示すことである。そのためには
$\text{WeakAss}_X(\mathcal{K})=\emptyset$ を示せば十分である。
Divisors on Spaces, Lemma
\ref{spaces-divisors-lemma-weakly-ass-zero} を参照せよ。
Divisors on Spaces, Lemma
\ref{spaces-divisors-lemma-ses-weakly-ass}
により
$\text{WeakAss}_X(\mathcal{K})\subset\text{WeakAss}_X(\mathcal{F})$
である。$\mathcal{F}$ は平坦なので、
Lemma \ref{spaces-flat-lemma-bourbaki-finite-type-general-base}
から
$\text{WeakAss}_X(\mathcal{F})\subset\text{Ass}_{X/Y}(\mathcal{F})$
を得る。仮定により、任意の点 $x$ で
$\text{Ass}_{X/Y}(\mathcal{F})$ に属するものは
ある $x'\in|X'|$ の像である。$g$ は平坦なので、局所環の写像
$\mathcal{O}_{X,\overline{x}}\to\mathcal{O}_{X',\overline{x}'}$
は忠実平坦であり、したがって写像
$$
\mathcal{F}_{\overline{x}}
\longrightarrow
(g^*\mathcal{F})_{\overline{x}'} =
\mathcal{F}_{\overline{x}} \otimes_{\mathcal{O}_{X, \overline{x}}}
\mathcal{O}_{X', \overline{x}'}
$$
は単射である
（Algebra, Lemma
\ref{algebra-lemma-faithfully-flat-universally-injective}
を参照せよ）。表示された矢印は
$\mathcal{F}_{\overline{x}}\to(g_*g^*\mathcal{F})_{\overline{x}}$
を経由するので、$\mathcal{K}_{\overline{x}}=0$ と結論できる。
したがって $x$ は $\mathcal{K}$ の弱随伴点ではあり得ず、証明が完了する。
\end{proof}




\section{平坦化関手}
\label{spaces-flat-section-F-zero}

\noindent
本節は More on Flatness, Section
\ref{flat-section-flattening-functors}
の類似である。初読時には本節を飛ばすことを勧める。

\begin{situation}
\label{spaces-flat-situation-iso}
$S$ をスキームとする。
$f:X\to B$ を $S$ 上の代数空間の射とする。
$u:\mathcal{F}\to\mathcal{G}$ を準連接
$\mathcal{O}_X$-加群の準同型とする。任意のスキーム $T$ で
$B$ 上にあるものに対して、
$u_T:\mathcal{F}_T\to\mathcal{G}_T$ により $u$ の $T$ への
基底変換を表す。言い換えれば、$u_T$ は $u$ の射影
$X_T=X\times_BT\to X$ による引き戻しである。
この状況で、関手
\begin{equation}
\label{spaces-flat-equation-iso}
F_{iso} : (\Sch/B)^{opp} \longrightarrow \textit{Sets}, \quad
T \longrightarrow \left\{
\begin{matrix}
\{*\} & \text{もし }u_T\text{ が同型ならば}, \\
\emptyset & \text{それ以外。}
\end{matrix}
\right.
\end{equation}
を考えられる。変種 $F_{inj}$、$F_{surj}$、$F_{zero}$ では、
$u_T$ がそれぞれ単射、全射、または零であることを要求する。
\end{situation}

\noindent
Situation \ref{spaces-flat-situation-iso} では、関手
$F_{iso}$、$F_{inj}$、$F_{surj}$、$F_{zero}$ を、関手
$(\Sch/S)^{opp}\to\textit{Sets}$ であって射
$F_{iso}\to B$、$F_{inj}\to B$、$F_{surj}\to B$、
$F_{zero}\to B$ を備えるものと考えることがある。
すなわち、$T$ が $S$ 上のスキームならば、
元 $h\in F_{iso}(T)$ は射 $h:T\to B$ であって、
$u$ の $h$ による基底変換が同型となるものである。
特に、$F_{iso}$ が代数空間であるというとき、それは対応する関手
$(\Sch/S)^{opp}\to\textit{Sets}$ が代数空間であるという意味である。

\begin{lemma}
\label{spaces-flat-lemma-iso-sheaf}
Situation \ref{spaces-flat-situation-iso} において、
関手 $F_{iso}$、$F_{inj}$、$F_{surj}$、$F_{zero}$ はいずれも
fpqc 位相に対する層条件を満たす。
\end{lemma}

\begin{proof}
$\{T_i\to T\}_{i\in I}$ を $B$ 上のスキームの fpqc 被覆とする。
% P09-ERR-1251
$X_i=X_{T_i}=X\times_ST_i$ および $u_i=u_{T_i}$ とおく。
$\{X_i\to X_T\}_{i\in I}$ は $X_T$ の fpqc 被覆であることに注意せよ。
Topologies on Spaces, Lemma
\ref{spaces-topologies-lemma-fpqc} を参照せよ。
特に、すべての $x\in|X_T|$ に対して、ある $i\in I$ と
$x_i\in|X_i|$ で $x$ に写るものが存在する。
局所環の写像
$\mathcal{O}_{X_T,\overline{x}}\to
\mathcal{O}_{X_i,\overline{x_i}}$
は平坦であり、したがって忠実平坦である
（Morphisms of Spaces, Section
\ref{spaces-morphisms-section-flat} を参照せよ）。
よって $(u_i)_{x_i}$ が単射、全射、全単射、または零であることと、
$(u_T)_x$ がそれぞれ単射、全射、全単射、または零であることは同値である。
補題が従う。
\end{proof}

\begin{lemma}
\label{spaces-flat-lemma-iso-go-up}
Situation \ref{spaces-flat-situation-iso} において、$X'\to X$ を代数空間の平坦射とする。
$u':\mathcal{F}'\to\mathcal{G}'$ により $u$ の $X'$ への引き戻しを表す。
$F'_{iso}$、$F'_{inj}$、$F'_{surj}$、$F'_{zero}$ により、
$\Sch/B$ 上の、$u'$ に付随する関手を表す。
\begin{enumerate}
\item $\mathcal{G}$ が有限型で、$|X'|\to|X|$ の像が
$\mathcal{G}$ の台を含むならば、$F_{surj}=F'_{surj}$ かつ
$F_{zero}=F'_{zero}$ である。
\item $\mathcal{F}$ が有限型で、$|X'|\to|X|$ の像が
$\mathcal{F}$ の台を含むならば、$F_{inj}=F'_{inj}$ かつ
$F_{zero}=F'_{zero}$ である。
\item $\mathcal{F}$ と $\mathcal{G}$ が有限型で、$|X'|\to|X|$ の像が
$\mathcal{F}$ と $\mathcal{G}$ の台を含むならば、
$F_{iso}=F'_{iso}$ である。
\end{enumerate}
\end{lemma}

\begin{proof}
準連接加群の写像 $v:\mathcal{H}\to\mathcal{E}$ を代数空間 $Y$ 上でとり、
$\varphi:Y'\to Y$ を代数空間の全射平坦射とする。このとき $v$ が
同型、単射、全射、または零であることと、$\varphi^*v$ がそれぞれ
同型、単射、全射、または零であることは同値である。
実際、すべての $y\in|Y|$ に対して $y'\in|Y'|$ が存在し、
局所環の写像
$\mathcal{O}_{Y,\overline{y}}\to\mathcal{O}_{Y',\overline{y'}}$
は忠実平坦である
（Morphisms of Spaces, Section
\ref{spaces-morphisms-section-flat} を参照せよ）。
もちろん、単射性または零であることを確認するには
$\mathcal{H}$ の台の点だけを見れば十分であり、全射性を確認するには
$\mathcal{E}$ の台の点だけを見れば十分である。さらに、補題の主張にある
有限型の仮定のもとでは、台をとる操作は基底変換と可換する。
Morphisms of Spaces, Lemma
\ref{spaces-morphisms-lemma-support-finite-type}
を参照せよ。したがって補題は明らかである。
\end{proof}

\noindent
有限型準連接加群のスキーム論的台は Morphisms of Spaces, Definition
\ref{spaces-morphisms-definition-scheme-theoretic-support}
で定義したことを思い出そう。

\begin{lemma}
\label{spaces-flat-lemma-iso-limits}
Situation \ref{spaces-flat-situation-iso} において考える。
\begin{enumerate}
\item $\mathcal{G}$ が有限型で、$\mathcal{G}$ のスキーム論的台が
$B$ 上準コンパクト
ならば、$F_{surj}$ は極限を保つ。
\item % P09-ERR-1252
$\mathcal{F}$ が有限型で、$\mathcal{F}$ のスキーム論的台が
$B$ 上準コンパクトならば、
$F_{zero}$ は極限を保つ。
\item $\mathcal{F}$ が有限型、$\mathcal{G}$ が有限表示で、
$\mathcal{F}$ と $\mathcal{G}$ のスキーム論的台が $B$ 上準コンパクトならば、
$F_{iso}$ は極限を保つ。
\end{enumerate}
\end{lemma}

\begin{proof}
(1) の証明。$i:Z\to X$ を $\mathcal{G}$ のスキーム論的台とし、
$\mathcal{G}$ を $Z$ 上の有限型準連接加群と考える。
$X$ を $Z$ で、$u$ を写像 $i^*\mathcal{F}\to\mathcal{G}$ で
置き換えてよい（詳細は省略する）。したがって $f$ が準コンパクトで
$\mathcal{G}$ が有限型であると仮定してよい。
$T=\lim_{i\in I}T_i$ をアフィン $B$-スキームの有向極限とし、
$u_T$ が全射であると仮定する。
$X_i=X_{T_i}=X\times_ST_i$ とおき、
$u_i=u_{T_i}:\mathcal{F}_i=\mathcal{F}_{T_i}
\to\mathcal{G}_i=\mathcal{G}_{T_i}$ とおく。
(1) を示すには、$u_i$ が全射となる $i$ が存在することを示せばよい。
$0\in I$ を選び、$I$ を $\{i\mid i\geq0\}$ で置き換える。
$f$ は準コンパクトなので $X_0$ は準コンパクトである。
したがって、全射エタール射 $\varphi_0:W_0\to X_0$ で
$W_0$ がアフィンスキームであるものを選べる。
$W=W_0\times_{T_0}T$ および
$W_i=W_0\times_{T_0}T_i$（$i\geq0$）とおく。
これらはアフィンスキームであり、全射エタール射
$\varphi:W\to X_T$ および $\varphi_i:W_i\to X_i$ を備える。
$W=\lim W_i$ であることに注意せよ。
したがって $\varphi^*u_T$ は全射であり、
$\varphi_i^*u_i$ が全射となる $i$ が存在することを示せば十分である。
よって問題はアフィンの場合、すなわち
Algebra, Lemma
\ref{algebra-lemma-module-map-property-in-colimit} の (2) に帰着された。

\medskip\noindent
(2) の証明。$\mathcal{F}$ は有限型で、そのスキーム論的台
% P09-ERR-1253
$Z\subset B$ は $B$ 上準コンパクトであると仮定する。
$T=\lim_{i\in I}T_i$ をアフィン $B$-スキームの有向極限とし、
$u_T$ が零であると仮定する。
$X_i=T_i\times_BX$ とおき、$u_i:\mathcal{F}_i\to\mathcal{G}_i$ により
引き戻しを表す。$0\in I$ を選び、$I$ を
$\{i\mid i\geq0\}$ で置き換える。$Z_0=Z\times_XX_0$ とおく。
Morphisms of Spaces, Lemma
\ref{spaces-morphisms-lemma-support-finite-type}
により、% P09-ERR-1254
$\mathcal{F}_i$ の台は $|Z_0|$ である。$|Z_0|$ は準コンパクトなので、
アフィンスキーム $W_0$ とエタール射 $W_0\to X_0$ で
$|Z_0|\subset\Im(|W_0|\to|X_0|)$ となるものを見つけられる。
$W=W_0\times_{T_0}T$ および
$W_i=W_0\times_{T_0}T_i$（$i\geq0$）とおく。
これらはアフィンスキームであり、エタール射
$\varphi:W\to X_T$ および $\varphi_i:W_i\to X_i$ を備える。
$W=\lim W_i$ であり、$\mathcal{F}_T$ と $\mathcal{F}_i$ の台は
それぞれ $|W|\to|X_T|$ と $|W_i|\to|X_i|$ の像に含まれる。
ここで % P09-ERR-1255
$\varphi^*u_T$ は単射であり、
$\varphi_i^*u_i$ が単射となる $i$ が存在することを示せば十分である。
よって問題はアフィンの場合、すなわち
Algebra, Lemma
\ref{algebra-lemma-module-map-property-in-colimit} の (1) に帰着された。

\medskip\noindent
(3) の証明。前二段落とまったく同様に、
Algebra, Lemma
\ref{algebra-lemma-module-map-property-in-colimit} の (3) を用いて証明できる。
次のように (1) と (2) から導くこともできる。
$T=\lim_{i\in I}T_i$ をアフィン $B$-スキームの有向極限とし、
$u_T$ が同型であると仮定する。(1) により、
$0\in I$ で $u_{T_0}$ が全射となるものが存在する。
$\mathcal{K}=\Ker(u_{T_0})$ とおき、準連接加群の写像
$v:\mathcal{K}\to\mathcal{F}_{T_0}$ を考える。
$i\geq0$ に対して、基底変換 $v_{T_i}$ が零であるための必要十分条件は
$u_i$ が同型であることである。さらに $v_T$ は零である。
$\mathcal{G}_{T_0}$ は有限表示、$\mathcal{F}_{T_0}$ は有限型であり、
$u_{T_0}$ は全射だから、$\mathcal{K}$ は有限型である
（Modules on Sites, Lemma
\ref{sites-modules-lemma-kernel-surjection-finite-onto-finite-presentation}）。
$\mathcal{K}$ の台が $\mathcal{F}_{T_0}$ の台に含まれ、
後者が $T_0$ 上準コンパクトであることは明らかである。
したがって (2) を適用して、$v_{T_i}$ が零となる $i$ が存在することが
分かる。
\end{proof}

\begin{lemma}
\label{spaces-flat-lemma-relate-zero-iso}
Situation \ref{spaces-flat-situation-iso} において、完全列
$$
\mathcal{F} \xrightarrow{u} \mathcal{G} \xrightarrow{v} \mathcal{H} \to 0
$$
が与えられているとする。このとき明らかな記法のもとで
% P09-ERR-1256
$F_{v,iso}=F_{u,zero}$ である。
\end{lemma}

\begin{proof}
引き戻しは右完全なので、
$$
\mathcal{F}_T\to\mathcal{G}_T\to\mathcal{H}_T\to0
$$
は完全である。これは任意のスキーム $T$ で $B$ 上にあるものに対して
成り立つ。したがって $u_T$ が全射であるための必要十分条件は
$v_T$ が同型であることである。
\end{proof}

\begin{lemma}
\label{spaces-flat-lemma-relate-zero-affine}
Situation \ref{spaces-flat-situation-iso} において、アフィン射 $i:Z\to X$ と
準連接 $\mathcal{O}_Z$-加群 $\mathcal{H}$ で
$\mathcal{G}=i_*\mathcal{H}$ となるものが与えられているとする。
$v:i^*\mathcal{F}\to\mathcal{H}$ を $u$ に随伴する写像とする。このとき
\begin{enumerate}
\item $F_{v,zero}=F_{u,zero}$ である。
\item $i$ が閉埋め込みならば $F_{v,surj}=F_{u,surj}$ である。
\end{enumerate}
\end{lemma}

\begin{proof}
$T$ を $B$ 上のスキームとする。$i_T:Z_T\to X_T$ により $i$ の
基底変換を表し、$\mathcal{H}_T$ により $\mathcal{H}$ の $Z_T$ への
引き戻しを表す。
$(i^*\mathcal{F})_T=i_T^*\mathcal{F}_T$ および
$i_{T,*}\mathcal{H}_T=(i_*\mathcal{H})_T$ であることに注意せよ。
第一の主張は引き戻しの可換性から、第二の主張は
Cohomology of Spaces, Lemma
\ref{spaces-cohomology-lemma-affine-base-change}
から従う。したがって $u_T$ と $v_T$ も随伴する写像である。
よって $u_T=0$ であるための必要十分条件は $v_T=0$ であることである。
これで (1) が証明された。(2) の場合、$u_T$ が全射であるための
必要十分条件は $v_T$ が全射であることである。実際、$u_T$ は
$$
\mathcal{F}_T \to
i_{T, *}i_T^*\mathcal{F}_T \xrightarrow{i_{T, *}v_T} i_{T, *}\mathcal{H}_T
$$
と分解し、$i_{T,*}$ は完全関手であって、$Z_T$ 上の準連接加群の圏を
準連接 $\mathcal{O}_{X_T}$-加群の圏へ充満忠実に埋め込む。
Morphisms of Spaces, Lemma
\ref{spaces-morphisms-lemma-i-star-equivalence}
を参照せよ。
\end{proof}

\begin{lemma}
\label{spaces-flat-lemma-relate-zero-affine-push}
Situation \ref{spaces-flat-situation-iso} において、アフィン射 $g:X\to X'$ が
与えられているとする。% P09-ERR-1257
$u'=f_*u:f_*\mathcal{F}\to f_*\mathcal{G}$ とおく。このとき
$F_{u,iso}=F_{u',iso}$、$F_{u,inj}=F_{u',inj}$、
$F_{u,surj}=F_{u',surj}$、かつ $F_{u,zero}=F_{u',zero}$ である。
\end{lemma}

\begin{proof}
Cohomology of Spaces, Lemma
\ref{spaces-cohomology-lemma-affine-base-change}
により $g_{T,*}u_T=u'_T$ である。さらに、% P09-ERR-1258
$g_{T,*}:\QCoh(\mathcal{O}_{X_T})\to\QCoh(\mathcal{O}_X)$
は忠実な完全関手であり、同型、単射、全射を反映する。
\end{proof}

\begin{situation}
\label{spaces-flat-situation-flat}
$S$ をスキームとする。
$f:X\to Y$ を $S$ 上の代数空間の射とする。
$\mathcal{F}$ を準連接 $\mathcal{O}_X$-加群とする。
任意のスキーム $T$ で $Y$ 上にあるものに対して、$\mathcal{F}_T$ により
$\mathcal{F}$ の $T$ への基底変換を表す。言い換えれば、
$\mathcal{F}_T$ は $\mathcal{F}$ の射影
$X_T=X\times_YT\to X$ による引き戻しである。
平坦加群の基底変換は平坦なので、関手
{\small
\begin{equation}
\label{spaces-flat-equation-flat}
F_{flat} : (\Sch/Y)^{opp} \longrightarrow \textit{Sets}, \quad
T \longrightarrow \left\{
\begin{matrix}
\{*\} & \text{もし }\mathcal{F}_T\text{ が次の基底上で平坦ならば：}T, \\
\emptyset & \text{それ以外。}
\end{matrix}
\right.
\end{equation}
}
を得る。
\end{situation}

\noindent
Situation \ref{spaces-flat-situation-flat} では、$F_{flat}$ を、関手
$(\Sch/S)^{opp}\to\textit{Sets}$ であって射
$F_{flat}\to Y$ を備えるものと考えることがある。
すなわち、$T$ が $S$ 上のスキームならば、
元 $h\in F_{flat}(T)$ は射 $h:T\to Y$ であって、
$\mathcal{F}$ の $h$ による基底変換が $T$ 上平坦となるものである。
特に、$F_{flat}$ が代数空間であるというとき、それは対応する関手
$(\Sch/S)^{opp}\to\textit{Sets}$ が代数空間であるという意味である。

\begin{lemma}
\label{spaces-flat-lemma-flat}
Situation \ref{spaces-flat-situation-flat} において考える。
\begin{enumerate}
\item 関手 $F_{flat}$ は fpqc 位相に対する層条件を満たす。
\item $f$ が準コンパクトかつ局所有限表示であり、
$\mathcal{F}$ が有限表示ならば、関手 $F_{flat}$ は極限を保つ。
\end{enumerate}
\end{lemma}

\begin{proof}
(1) は次の主張から従う。$T'\to T$ が $Y$ 上の代数空間の全射平坦射ならば、
$\mathcal{F}_{T'}$ が $T'$ 上平坦であるための必要十分条件は
$\mathcal{F}_T$ が $T$ 上平坦であることである。
Morphisms of Spaces, Lemma
\ref{spaces-morphisms-lemma-base-change-module-flat}
を参照せよ。(2) は、$f$ がさらに準分離
（すなわち $f$ が有限表示）ならば
Limits of Spaces, Lemma
\ref{spaces-limits-lemma-descend-flat}
から従う。一般の場合は、まず基底がアフィンである場合に帰着し、
次に $X$ を有限個のアフィンで被覆して準分離の場合に帰着する。
詳細は省略する。
\end{proof}

\section{写像を零にすること}
\label{spaces-flat-section-zero-map}

\noindent
本節には、対応するスキームの章に類似する節はない。

\begin{situation}
\label{spaces-flat-situation-somewhat-closed}
$S=\Spec(R)$ をアフィンスキームとする。$X$ を $S$ 上の代数空間とする。
$u:\mathcal{F}\to\mathcal{G}$ を準連接 $\mathcal{O}_X$-加群の写像とする。
% P09-ERR-1259
$\mathcal{G}$ は $S$ 上平坦であると仮定する。
\end{situation}

\begin{lemma}
\label{spaces-flat-lemma-F-zero-somewhat-closed}
Situation \ref{spaces-flat-situation-somewhat-closed} において考える。
$T\to S$ をスキームの準コンパクト射で、基底変換 $u_T$ が零となるものとする。
このとき、% P09-ERR-1260
閉部分スキーム $Z\subset S$ で、(a) $T\to S$ が $Z$ を経由し、
(b) 基底変換 $u_Z$ が零となるものが存在する。
$\mathcal{F}$ が有限型 $\mathcal{O}_X$-加群で、その $\mathcal{F}$ のスキーム論的台が
準コンパクトならば、$Z\to S$ を有限表示にとれる。
\end{lemma}

\begin{proof}
$U\to X$ を全射エタール射で、$U=\coprod U_i$ がアフィンスキームの
非交和となるものとする
（Properties of Spaces, Lemma
\ref{spaces-properties-lemma-cover-by-union-affines} を参照せよ）。
Lemma \ref{spaces-flat-lemma-iso-go-up} により、$X$ を $U$ で置き換えてよい。
言い換えれば、$X=\coprod X_i$ がアフィンスキーム $X_i$ の非交和であると
仮定してよい。補題を $u_i=u|_{X_i}$ に対して証明できると仮定する。
すると、閉部分スキーム $Z_i\subset S$ で、$T\to S$ が $Z_i$ を経由し、
$u_{i,Z_i}$ が零となるものを得る。
$Z_i=\Spec(R/I_i)\subset\Spec(R)=S$ ならば、
$Z=\Spec(R/\sum I_i)$ ととればよい。したがって
$X=\Spec(A)$ がアフィンであると仮定してよい。

\medskip\noindent
有限アフィン開被覆 $T=T_1\cup\ldots\cup T_m$ を選ぶ。
$T$ を $\coprod_{j=1,\ldots,m}T_j$ で置き換えてよいことは明らかである。
したがって $T$ がアフィンであると仮定してよい。
$T=\Spec(R')$ と書く。% P09-ERR-1261
$u:M\to N$ を $A$-加群の準同型とし、
$u:\mathcal{F}\to\mathcal{G}$ に対応するものとする。$N$ は平坦な $R$-加群である。
これは $\mathcal{G}$ が $S$ 上平坦であることによる。補題の仮定は、合成
$$
M \otimes_R R' \to N \otimes_R R'
$$
が零であることを意味する。$z\in M$ とする。
Lazard の定理
（Algebra, Theorem \ref{algebra-theorem-lazard}）
および $\otimes$ が余極限と可換することにより、
自由 $R$-加群 $F_z$、元 $\tilde z\in F_z$、写像 $F_z\to N$ で、
$u(z)$ が $\tilde z$ の像となり、$\tilde z$ が
$F_z\otimes_RR'$ において零へ写るものを見つけられる。
基底 $\{e_{z,\alpha}\}$ を $F_z$ に選び、
$\tilde z=\sum f_{z,\alpha}e_{z,\alpha}$ と書く。ここで
$f_{z,\alpha}\in R$ である。$I\subset R$ を、元 $f_{z,\alpha}$（$z$ が
$M$ のすべての元を動くときのもの）で生成されるイデアルとする。
構成により $I$ は $R'$ において零へ写り、元 $\tilde z$ は
$F_z/IF_z$ において、したがって $N/IN$ において零へ写る。
よって $Z=\Spec(R/I)$ はこの場合の問題の解である。

\medskip\noindent
$\mathcal{F}$ は有限型で準コンパクトなスキーム論的台をもつと仮定する。
$Z=\Spec(R/I)$ と書く。$I=\bigcup I_\lambda$ を有限生成イデアルの
フィルター付き和として書く。$Z_\lambda=\Spec(R/I_\lambda)$ とおくと、
$Z=\colim Z_\lambda$ である。$u_Z$ は零なので、$u_{Z_\lambda}$ は、
ある $\lambda$ に対して零である（Lemma \ref{spaces-flat-lemma-iso-limits} による）。
これで補題の証明は完了する。
\end{proof}

\begin{lemma}
\label{spaces-flat-lemma-F-zero-module-map}
$A$ を環とする。$u:M\to N$ を $A$-加群の写像とする。
$N$ が射影 $A$-加群ならば、イデアル $I\subset A$ で、任意の環写像
$\varphi:A\to B$ に対して次が同値となるものが存在する。
\begin{enumerate}
\item $u\otimes1:M\otimes_AB\to N\otimes_AB$ は零である。
\item $\varphi(I)=0$ である。
\end{enumerate}
\end{lemma}

\begin{proof}
$N$ は射影なので、射影 $A$-加群 $C$ で、% P09-ERR-1262
$F=N\oplus C$ が自由 $R$-加群となるものを見つけられる。
% P09-ERR-1263
$u$ を $u\oplus1:F=M\oplus C\to N\oplus C$ で置き換えることにより、
$N$ が自由であると仮定してよい。この場合、$I$ を、$A$ のイデアルで、
$\Im(u)$ のすべての元を $N$ の固定した基底で表した係数により生成される
ものとする。
\end{proof}

\begin{lemma}
\label{spaces-flat-lemma-F-zero-somewhat-closed-points}
Situation \ref{spaces-flat-situation-somewhat-closed} において考える。
$T\subset S$ を部分集合とし、$s\in S$ は $T$ の閉包に属するとする。
$t\in T$ に対して、$u_t$ を $u$ の $X_t$ への引き戻しとし、
$u_s$ を $u$ の $X_s$ への引き戻しとする。
$X$ が $S$ 上局所有限表示であり、
$\mathcal{G}$ が有限表示\footnote{$X$ が $S$ 上局所有限型で、
$\mathcal{G}$ が $S$ に関して有限表示であると仮定するだけで十分であるが、
この概念は代数空間の枠組みではまだ定義されていない。
スキームに対する定義は More on Morphisms, Section
\ref{more-morphisms-section-finite-type-finite-presentation}
で与えられている。}であり、$u_t=0$ がすべての $t\in T$ に対して
成り立つならば、$u_s=0$ である。
\end{lemma}

\begin{proof}
$u_s$ が零であるかどうかの確認は、ファイバー $X_s$ 上エタール局所的である。
したがって点 $x\in|X_s|\subset|X|$ を選び、エタール近傍で確認してよい。
Proposition \ref{spaces-flat-proposition-finite-presentation-flat-at-point}
にあるような
$$
\xymatrix{
(X, x) \ar[d] & (X', x') \ar[l]^g \ar[d] \\
(S, s) & (S', s') \ar[l]
}
$$
を選ぶ。$T'\subset S'$ を $T$ の逆像とする。
$s'$ は $T'$ の閉包に属する。これは $S'\to S$ が開だからである。
したがって、次の段落で記述する代数の問題に帰着される。

\medskip\noindent
$R$-加群の写像 $u:M\to N$ で、$N$ が射影 $R$-加群であり、
$u_t:M\otimes_R\kappa(t)\to N\otimes_R\kappa(t)$
が各 $t\in T$ に対して零となるものがある。問題は $u_s=0$ を示すことである。
$I\subset R$ を Lemma \ref{spaces-flat-lemma-F-zero-module-map} で定義した
イデアルとする。すると $I$ は $\kappa(t)$ において、すべての $t\in T$ に対して
零へ写る。したがって $T\subset V(I)$ である。
$s$ は $T$ の閉包に属するので $s\in V(I)$ である。
よって $u_s=0$ である。
\end{proof}

\noindent
$f:X\to B$ が射影射で、$B$ がネーター・スキームである場合に、補題
\ref{spaces-flat-lemma-F-zero-closed-pure} または
\ref{spaces-flat-lemma-F-zero-closed-proper} のいずれかを
``単純な''直接証明で示し、補題
\ref{spaces-flat-lemma-F-zero-somewhat-closed} および
\ref{spaces-flat-lemma-F-zero-somewhat-closed-points} の証明で用いたような議論を
使えれば興味深い。古典的な証明は次のようになる。(a) 相対的に豊富な可逆層
$\mathcal{O}_X(1)$ を選び、(b)
$u_n:f_*\mathcal{F}(n)\to f_*\mathcal{G}(n)$ とおき、(c)
$f_*\mathcal{G}(n)$ がすべての $n\gg 0$ に対して有限局所自由層であることを
観察し、(d) $F_{zero}$ が $u_n$ のある $n\gg 0$ における零点軌跡で
表されることを示す。

\begin{lemma}
\label{spaces-flat-lemma-F-zero-closed-pure}
Situation \ref{spaces-flat-situation-iso} において考える。次を仮定する。
\begin{enumerate}
\item $f$ は有限表示である。
\item $\mathcal{G}$ は $B$ 上平坦かつ $B$ に関して純粋な有限表示加群である。
\end{enumerate}
このとき $F_{zero}$ は代数空間であり、$F_{zero}\to B$ は閉埋め込みである。
$\mathcal{F}$ が有限型ならば、$F_{zero}\to B$ は有限表示である。
\end{lemma}

\begin{proof}
Lemma \ref{spaces-flat-lemma-finite-type-flat-pure-is-universal} により、加群
$\mathcal{G}$ は $B$ に関して普遍的に純粋である。$F_{zero}$ が代数空間で
あることを示すには、$F_{zero}\to B$ が表現可能であることを示せば十分で
ある（Spaces, Lemma \ref{spaces-lemma-representable-over-space} を参照）。
$B'\to B$ という射を取り、$B'$ はスキームとする。$u':\mathcal{F}'\to\mathcal{G}'$
を $u$ の引き戻しとして定め、$X'=X_{B'}$ とおく。このとき対応する関手
$F'_{zero}$ は $F_{zero}\times_B B'$ に等しい。したがって $B$ がスキームで
ある場合に帰着する。

\medskip\noindent
$B$ はスキームであると仮定する。$F_{zero}$ が $B$ の閉部分スキームによって
表現可能であることを示す。Lemma \ref{spaces-flat-lemma-iso-sheaf} と Descent の補題
\ref{descent-lemma-closed-immersion} および
\ref{descent-lemma-descent-data-sheaves} により、この問題は $B$ の
エタール位相に関して局所的である。$b\in B$ とする。まず $B$ を $b$ の
アフィン近傍で置き換える。
図式
$$
\xymatrix{
X \ar[d] & X' \ar[l]^g \ar[d] \\
B & B' \ar[l]
}
$$
と、$b'\in B'$ を取り、$b\in B$ に写るものとする。Lemma
\ref{spaces-flat-lemma-finite-presentation-flat-along-fibre} のように選ぶ。エタール局所で
作業しているので、$B$ を $B'$ で置き換え、図式
$$
\xymatrix{
X \ar[rd] & & X' \ar[ll]^g \ar[ld] \\
& B
}
$$
であって、$B$ と $X'$ がアフィンで、$\Gamma(X',g^*\mathcal{G})$ が
$\Gamma(B,\mathcal{O}_B)$-加群として射影的であり、かつ
$g(|X'|)\supset|X_b|$ であるものを仮定してよい。開部分空間 $U\subset X$ を取り、
$|U|=g(|X'|)$ とする。Divisors on Spaces, Lemma
\ref{spaces-divisors-lemma-relative-assassin-constructible} により、
$$
E=\{t\in B:\text{Ass}_{X_t}(\mathcal{G}_t)\subset|U_t|\}=
\{t\in B:\text{Ass}_{X/B}(\mathcal{G})\cap|X_t|\subset|U_t|\}
$$
は $B$ において構成可能である。Lemma \ref{spaces-flat-lemma-pure-on-top} の (2) により、
$E$ は $\Spec(\mathcal{O}_{B,b})$ を含む。Morphisms, Lemma
\ref{morphisms-lemma-constructible-containing-open} により、$E$ は $b$ の
開近傍を含む。したがって $B$ を $b$ のより小さいアフィン近傍で置き換えれば、
$$
\text{Ass}_{X/B}(\mathcal{G})\subset g(|X'|)
$$
と仮定できる。

\medskip\noindent
Lemma \ref{spaces-flat-lemma-universally-separating} から、% P09-ERR-1264
$u:\mathcal{F}\to\mathcal{G}$ が単射であることと
$g^*u:g^*\mathcal{F}\to g^*\mathcal{G}$ が単射であることは同値であり、
これは任意の基底変換後にも成り立つ。したがって、定理の仮定に加えて
$X\to B$ がアフィンスキームの射であり、
$\Gamma(X,\mathcal{G})$ が $\Gamma(B,\mathcal{O}_B)$-加群として射影的である
場合に帰着した。この場合は Lemma \ref{spaces-flat-lemma-F-zero-module-map} から直ちに
従う。

\medskip\noindent
残っているのは、$F_{zero}\to B$ が有限表示であることを、
$\mathcal{F}$ が有限型の場合に示すことである。これは Lemma \ref{spaces-flat-lemma-iso-limits} と Limits of
Spaces, Proposition
\ref{spaces-limits-proposition-characterize-locally-finite-presentation} を
組み合わせれば従う。
\end{proof}

\begin{lemma}
\label{spaces-flat-lemma-F-zero-closed-proper}
Situation \ref{spaces-flat-situation-iso} において考える。次を仮定する。
\begin{enumerate}
\item $f$ は局所有限表示である。
\item $\mathcal{G}$ は有限表示の $\mathcal{O}_X$-加群で、$B$ 上平坦である。
\item $\mathcal{G}$ の台は $B$ 上固有である。
\end{enumerate}
このとき関手 $F_{zero}$ は代数空間であり、$F_{zero}\to B$ は閉埋め込みである。
$\mathcal{F}$ が有限型ならば、$F_{zero}\to B$ は有限表示である。
\end{lemma}

\begin{proof}
$f$ が有限表示ならば、これは補題
\ref{spaces-flat-lemma-F-zero-closed-pure} および \ref{spaces-flat-lemma-proper-pure} から直ちに
従う。この場合だけが重要なので、この補題の仮定が許す、$f$ が準分離でも
準コンパクトでもない可能性を扱う残りの証明は読み飛ばすことを勧める。

\medskip\noindent
$i:Z\to X$ を、$\mathcal{G}$ の第 0 フィッティング・イデアルで切り取られる
閉部分空間として取る（Divisors on Spaces, Section
\ref{spaces-divisors-section-fitting-ideals}）。仮定により $Z\to B$ は固有で
ある（Derived Categories of Spaces, Section
\ref{spaces-perfect-section-proper-over-base} を参照）。一方、$i$ は有限表示
である（Divisors on Spaces, Lemma
\ref{spaces-divisors-lemma-fitting-ideal-of-finitely-presented} および
Morphisms of Spaces, Lemma
\ref{spaces-morphisms-lemma-closed-immersion-finite-presentation}）。
$\mathcal{O}_Z$-加群 $\mathcal{H}$ が有限型の準連接加群として存在し、
$i_*\mathcal{H}=\mathcal{G}$ となる（Divisors on Spaces, Lemma
\ref{spaces-divisors-lemma-on-subscheme-cut-out-by-Fit-0}）。実際、Algebra,
Lemma \ref{algebra-lemma-finitely-presented-over-subring} により
$\mathcal{H}$ は $\mathcal{O}_Z$-加群として有限表示である（詳細は省略する）。
すると $F_{zero}$ は、$F_{zero}$、すなわち
$i^*\mathcal{F}\to\mathcal{H}$（$u$ に随伴する写像）に対する関手、と同じである
（Lemma \ref{spaces-flat-lemma-relate-zero-affine} を参照）。$\mathcal{H}$ は
$B$ に関して平坦である。これは $\mathcal{G}$ も同様だからである（茎上で確認せよ。
詳細は省略）。
さらに、$\mathcal{F}$ が有限型ならば $i^*\mathcal{F}$ も有限型であることに
注意せよ。したがってこの補題は証明の最初の段落で扱った場合に帰着される。
\end{proof}


\section{写像の平坦化}
\label{spaces-flat-section-flattening-map}

\noindent
この節は More on Flatness, Section \ref{flat-section-flattening-map} に対応する。
特に、以下の結果は More on Flatness, Theorem
\ref{flat-theorem-flattening-map} の変種である。

\begin{theorem}
\label{spaces-flat-theorem-flattening-map}
Situation \ref{spaces-flat-situation-iso} において、次を仮定する。
\begin{enumerate}
\item $f$ は有限表示である。
\item $\mathcal{F}$ は $B$ 上平坦かつ $B$ に関して純粋な有限表示加群であり、
\item $u$ は全射である。
\end{enumerate}
このとき $F_{iso}$ は閉埋め込み $Z\to B$ によって表現される。さらに、
$Z\to S$ は、$\mathcal{G}$ が有限表示ならば有限表示である。
\end{theorem}

\begin{proof}
$\mathcal{K}=\Ker(u)$ とおき、包含写像を $v:\mathcal{K}\to\mathcal{F}$ とする。
Lemma \ref{spaces-flat-lemma-relate-zero-iso} により
$F_{u,iso}=F_{v,zero}$ である。$v$ に Lemma
\ref{spaces-flat-lemma-F-zero-closed-pure} を適用すると、
$F_{u,iso}=F_{v,zero}$ は $B$ の閉部分空間によって表現される。
$\mathcal{K}$ は有限型であることに注意せよ。これは $\mathcal{G}$ が有限表示ならば
成り立つ（Modules on Sites, Lemma
\ref{sites-modules-lemma-kernel-surjection-finite-onto-finite-presentation} による。
したがって補題の最後の主張も得られる。
\end{proof}

\begin{lemma}
\label{spaces-flat-lemma-F-iso-closed}
Situation \ref{spaces-flat-situation-iso} において、次を仮定する。
\begin{enumerate}
\item $f$ は局所有限表示である。
\item $\mathcal{F}$ は局所有限表示で $B$ 上平坦である。
\item $\mathcal{F}$ の台は $B$ 上固有であり、
\item $u$ は全射である。
\end{enumerate}
このとき関手 $F_{iso}$ は代数空間であり、$F_{iso}\to B$ は閉埋め込みである。
$\mathcal{G}$ が有限表示ならば、$F_{iso}\to B$ は有限表示である。
\end{lemma}

\begin{proof}
$\mathcal{K}=\Ker(u)$ とおき、包含写像を $v:\mathcal{K}\to\mathcal{F}$ とする。
Lemma \ref{spaces-flat-lemma-relate-zero-iso} により
$F_{u,iso}=F_{v,zero}$ である。$v$ に Lemma
\ref{spaces-flat-lemma-F-zero-closed-proper} を適用すると、
$F_{u,iso}=F_{v,zero}$ は $B$ の閉部分空間によって表現される。
$\mathcal{K}$ は有限型であることに注意せよ。これは $\mathcal{G}$ が有限表示ならば
成り立つ（Modules on Sites, Lemma
\ref{sites-modules-lemma-kernel-surjection-finite-onto-finite-presentation} による。
したがって補題の最後の主張も得られる。
\end{proof}

\noindent
Quot 関手を論じる際に、次の（容易な）結果を用いる。

\begin{lemma}
\label{spaces-flat-lemma-F-surj-open}
Situation \ref{spaces-flat-situation-iso} において、次を仮定する。
\begin{enumerate}
\item $f$ は局所有限表示である。
\item $\mathcal{G}$ は有限型である。
\item $\mathcal{G}$ の台は $B$ 上固有である。
\end{enumerate}
このとき $F_{surj}$ は代数空間であり、$F_{surj}\to B$ は開埋め込みである。
\end{lemma}

\begin{proof}
$\Coker(u)$ を考える。$\Coker(u_T)=\Coker(u)_T$ は任意の $T/B$ に対して
成り立つことに注意せよ。有限型準連接加群の台の形成は
基底変換と可換する（Morphisms of Spaces, Lemma
\ref{spaces-morphisms-lemma-support-covering}）。したがって $F_{surj}$ は、
$B$ の開部分空間で、次の開集合に対応するものによって表現される。
$$
|B|\setminus|f|(\text{Supp}(\Coker(u)))
$$
（Properties of Spaces, Lemma
\ref{spaces-properties-lemma-open-subspaces} を参照）。これは開集合である。
なぜなら $|f|$ は $\text{Supp}(\mathcal{G})$ 上閉であり、
$\text{Supp}(\Coker(u))$ は $\text{Supp}(\mathcal{G})$ の閉部分集合だからである。
\end{proof}

\section{局所の場合の平坦化}
\label{spaces-flat-section-flattening-local}

\noindent
この節は More on Flatness, Section \ref{flat-section-flattening-local} に対応する。

\begin{lemma}
\label{spaces-flat-lemma-freebie}
$S$ を閉点 $s$ をもつヘンゼル局所環のスペクトルとする。$X\to S$ を
局所有限型の代数空間の射とする。$\mathcal{F}$ を有限型の準連接
$\mathcal{O}_X$-加群とする。$E\subset|X_s|$ を部分集合とする。次の性質を
もつ閉部分スキーム $Z\subset S$ が存在する。任意の点付きスキームの射
$(T,t)\to(S,s)$ に対して、次の条件は同値である。
\begin{enumerate}
\item $\mathcal{F}_T$ は $T$ 上、$|X_t|$ のうち
$E\subset|X_s|$ の点に写るすべての点で平坦である。
\item $\Spec(\mathcal{O}_{T,t})\to S$ は $Z$ を経由する。
\end{enumerate}
さらに、$X\to S$ が局所有限表示で、$\mathcal{F}$ が有限表示であり、
$E\subset|X_s|$ が閉かつ準コンパクトならば、$Z\to S$ は有限表示である。
\end{lemma}

\begin{proof}
スキーム $U$ とエタール射 $\varphi:U\to X$ を選ぶ。$E'\subset|U_s|$ を
$E$ の逆像とする。$E'\to E$ が全射ならば、条件 (1) は次と同値である。
$(\varphi^*\mathcal{F})_T$ が $T$ 上、$|U_t|$ のうち
$E'\subset|U_t|$ の点に写るすべての点で平坦である。$\varphi$ を全射となるように選べば、
More on Flatness, Lemma \ref{flat-lemma-freebie} によるスキームの場合へ帰着する。
$E$ が閉かつ準コンパクトならば、アフィン $U$ を選び、$E'\to E$ が全射となる
ようにできる。このとき $E'$ は閉かつ準コンパクトであり、最後の主張は
More on Flatness, Lemma \ref{flat-lemma-freebie} の最後の主張から従う。
\end{proof}

\section{普遍平坦化}
\label{spaces-flat-section-flattening-final}

\noindent
この節は More on Flatness, Section \ref{flat-section-flattening-final} に対応する。
主な目的は Lemma \ref{spaces-flat-lemma-when-universal-flattening} を証明することである。
しかし、この結果をスキームの場合の対応する結果から直接導く方法は見いだせない。
したがって、ここでいくつかの材料を再構築する必要がある。まず定義から始める。

\begin{definition}
\label{spaces-flat-definition-flattening}
スキーム $S$ とする。代数空間の射 $X\to Y$ を取り、これは $S$ 上とする。
$\mathcal{F}$ を準連接 $\mathcal{O}_X$-加群とする。「$\mathcal{F}$ の普遍平坦化が
存在する」とは、Situation \ref{spaces-flat-situation-flat} で定義した関手 $F_{flat}$ が
代数空間であることを言う。
$X$ の普遍平坦化が存在するとは、$\mathcal{O}_X$ の普遍平坦化が存在することを言う。
\end{definition}

\noindent
これは少し不満足な定義である。というのも、代数空間の任意の単射が表現可能かどうか
（More on Morphisms of Spaces, Section
\ref{spaces-more-morphisms-section-monomorphisms}）を知らないため、ここでの
普遍平坦化の定義はスキームの場合に用いたものと一致しないからである。このために
混乱が生じないことを願う。

\begin{lemma}
\label{spaces-flat-lemma-pre-flat-dimension-n}
スキーム $S$ とする。$f:X\to Y$ を局所有限型の代数空間の射とする。
$\mathcal{F}$ を有限型の準連接 $\mathcal{O}_X$-加群とする。$n\geq0$ とする。
次の条件は同値である。
\begin{enumerate}
\item ある可換図式
$$
\xymatrix{
U \ar[d]_\varphi \ar[r] & V \ar[d] \\
X \ar[r] & Y
}
$$
で、上下の射が全射エタール、$U$ と $V$ がスキームであるものに対して、層
$\varphi^*\mathcal{F}$ は $V$ 上、次元 $\geq n$ で平坦である（More on Flatness,
Definition \ref{flat-definition-flat-dimension-n}）。
\item 任意の可換図式
$$
\xymatrix{
U \ar[d]_\varphi \ar[r] & V \ar[d] \\
X \ar[r] & Y
}
$$
で、上下の射がエタール、$U$ と $V$ がスキームであるものに対して、層
$\varphi^*\mathcal{F}$ は $V$ 上、次元 $\geq n$ で平坦である。
\item $x\in|X|$ で $\mathcal{F}$ が $x$ において $Y$ 上平坦でないものに対し、
$x/f(x)$ の超越次数は $<n$ である（Morphisms of Spaces, Definition
\ref{spaces-morphisms-definition-dimension-fibre}）。
\end{enumerate}
これが成り立つならば、任意の基底変換 $Y'\to Y$ の後にも成り立つ。
\end{lemma}

\begin{proof}
(1) のような図式があるとする。More on Flatness, Lemma
\ref{flat-lemma-pre-flat-dimension-n} における条件の同値性により、(1) と (3) は
同値である。しかし (3) は $\varphi^*\mathcal{F}$ に継承される。ここで
$U\to V$ は (2) のような任意の射である。ゆえに、再びスキームの場合の結果により、
(3) は (2) を含意する。スキームの場合の結果は基底変換に関する主張も含意する。
\end{proof}

\begin{definition}
\label{spaces-flat-definition-flat-dimension-n}
スキーム $S$ とする。局所有限型代数空間の射 $f:X\to Y$ を取り、これは $S$ 上とする。
$\mathcal{F}$ を有限型の準連接 $\mathcal{O}_X$-加群とする。$n\geq0$ とする。
Lemma \ref{spaces-flat-lemma-pre-flat-dimension-n} の同値な条件が満たされるとき、
「$\mathcal{F}$ は $Y$ 上、次元 $\geq n$ で平坦である」と言う。
\end{definition}

\begin{situation}
\label{spaces-flat-situation-flat-dimension-n}
スキーム $S$ とする。$f:X\to Y$ を代数空間の射で、$S$ 上局所有限型なものと
する。$\mathcal{F}$ を有限型の準連接 $\mathcal{O}_X$-加群とする。
任意のスキーム $T$ で $Y$ 上にあるものに対して、$\mathcal{F}_T$ により
$\mathcal{F}$ の $T$ への基底変換を表す。言い換えれば、$\mathcal{F}_T$ は
$\mathcal{F}$ の射影 $X_T=X\times_YT\to X$ による引き戻しである。
$f_T:X_T\to T$ は有限型であり、$\mathcal{F}_T$ は有限型
$\mathcal{O}_{X_T}$-加群であることに注意せよ（Morphisms of Spaces, Lemma
\ref{spaces-morphisms-lemma-base-change-finite-type} および Modules on Sites,
Lemma \ref{sites-modules-lemma-local-pullback}）。$n\geq0$ とする。
Definition \ref{spaces-flat-definition-flat-dimension-n} と Lemma
\ref{spaces-flat-lemma-pre-flat-dimension-n} により、関手
\begin{equation}
\label{spaces-flat-equation-flat-dimension-n}
F_n : (\Sch/Y)^{opp} \longrightarrow \textit{Sets}, \quad
T \longrightarrow \left\{
\begin{matrix}
\{*\} & \text{if }\mathcal{F}_T\text{ is flat over }T\text{ in }\dim \geq n, \\
\emptyset & \text{else.}
\end{matrix}
\right.
\end{equation}
を得る。
\end{situation}

\noindent
Situation \ref{spaces-flat-situation-flat-dimension-n} では、$F_n$ を、関手
$(\Sch/S)^{opp}\to\textit{Sets}$ であって射 $F_n\to Y$ を備えるものと
考えることがある。すなわち、$T$ が $S$ 上のスキームならば、
$h\in F_n(T)$ は射 $h:T\to Y$ であって、$\mathcal{F}$ の $h$ による
基底変換が $T$ 上、次元 $\dim\geq n$ で平坦となるものである。特に、$F_n$ が
代数空間であるというとき、それは対応する関手
$(\Sch/S)^{opp}\to\textit{Sets}$ が代数空間であるという意味である。

\begin{lemma}
\label{spaces-flat-lemma-flat-dimension-n}
Situation \ref{spaces-flat-situation-flat-dimension-n} において考える。
\begin{enumerate}
\item 関手 $F_n$ は fpqc 位相に対する層条件を満たす。
\item $f$ が準コンパクトかつ局所有限表示で、$\mathcal{F}$ が有限表示ならば、
関手 $F_n$ は極限を保つ。
\end{enumerate}
\end{lemma}

\begin{proof}
(1) の証明。$\{T_i\to T\}$ を fpqc 被覆とし、$T$ は $Y$ 上のスキームとする。
$F_n(T_i)$ がすべての $i$ に対して空でないならば $F_n(T)$ が空でないことを
示さなければならない。Lemma \ref{spaces-flat-lemma-pre-flat-dimension-n} の (1) のような
図式を選ぶ。$F'_n$ を、$\varphi^*\mathcal{F}$ と射 $U\to V$ に対応する
関手とする。More on Flatness, Lemma \ref{flat-lemma-flat-dimension-n} により、
$F'_n$ には層条件が成り立つ。したがって $F_n$ にも層条件が成り立つ。
実際、$T\to Y$ に対して Lemma \ref{spaces-flat-lemma-pre-flat-dimension-n} により
$F_n(T)=F'_n(V\times_YT)$ であり、
$\{V\times_YT_i\to V\times_YT\}$ は fpqc 被覆だからである。

\medskip\noindent
(2) の証明。$T=\lim_{i\in I}T_i$ とし、$T_i$ を $Y$ 上のアフィンスキームとする。
これはフィルター付き極限であり、$F_n(T)$ が空でないと仮定する。
$F_n(T_i)$ がある $i$ に対して空でないことを示さなければならない。Lemma
\ref{spaces-flat-lemma-pre-flat-dimension-n} の (1) のような図式を選ぶ。$i\in I$ を固定し、
アフィン開 $W_i\subset V\times_YT_i$ で $T_i$ 上に全射的に写るものを選ぶ。
$i'\geq i$ に対して、$W_{i'}$ を $W_i$ の $V\times_YT_{i'}$ における逆像とし、
$W\subset V\times_YT$ を $W_i$ の逆像とする。このとき
% P09-ERR-1270
$W=\lim_{i'\geq i}W_i$ は $V$ 上のアフィンスキームのフィルター付き極限で
ある。再び Lemma \ref{spaces-flat-lemma-pre-flat-dimension-n} により、$F'_n(W_{i'})$ が
ある $i'\geq i$ に対して空でないことを示せば十分である。しかし
$F'_n(W)$ は空でない。これは $F_n(T)=F'_n(V\times_YT)$ が空でないという
仮定による。
したがって More on Flatness, Lemma \ref{flat-lemma-flat-dimension-n} を適用して
結論を得る。
\end{proof}

\begin{lemma}
\label{spaces-flat-lemma-localize-flat-dimension-n}
Situation \ref{spaces-flat-situation-flat-dimension-n} において考える。
$h:X'\to X$ をエタール射とする。$\mathcal{F}'=h^*\mathcal{F}$ および
$f'=f\circ h$ とおく。$F_n'$ を $(f':X'\to Y,\mathcal{F}')$ に対応する
(\ref{spaces-flat-equation-flat-dimension-n}) とする。このとき $F_n$ は $F_n'$ の
部分関手であり、$h(X')\supset\text{Ass}_{X/Y}(\mathcal{F})$ ならば
$F_n=F'_n$ である。
\end{lemma}

\begin{proof}
Lemma \ref{spaces-flat-lemma-pre-flat-dimension-n} の (1) のような $U\to X$、$V\to Y$、
$U\to V$ を選ぶ。全射エタール射 $U'\to U\times_XX'$ で $U'$ がスキームと
なるものを選ぶ。このとき、二つの関手 $F_{U,n}$ と $F_{U',n}$（$U'\to U$ と
$\mathcal{F}|_U$ によって $V$ 上定まるもの）に対して
補題の主張が成り立つ（More on Flatness, Lemma
\ref{flat-lemma-localize-flat-dimension-n} を参照）。一方、Lemma
\ref{spaces-flat-lemma-pre-flat-dimension-n} により、与えられた $T\to Y$ に対して
$F_n(T)=F_{U,n}(V\times_YT)$ および
$F'_n(T)=F_{U',n}(V\times_YT)$ である。これで補題は証明された。
\end{proof}

\begin{theorem}
\label{spaces-flat-theorem-flat-dimension-n-representable}
Situation \ref{spaces-flat-situation-flat-dimension-n} において考える。さらに、$f$ は有限表示、
$\mathcal{F}$ は有限表示 $\mathcal{O}_X$-加群で、$\mathcal{F}$ は $Y$ に関して
純粋であると仮定する。このとき $F_n$ は代数空間であり、$F_n\to Y$ は
有限表示の単射である。
\end{theorem}

\begin{proof}
Lemma \ref{spaces-flat-lemma-flat-dimension-n} により、関手 $F_n$ は fppf 位相に対する層で
ある。$F_n\to Y$ は $(\Sch/S)_{fppf}$ 上の層の単射なので、
$\Delta:F_n\to F_n\times F_n$ は対角射
$\Delta_Y:Y\to Y\times_SY$ の引き戻しである。したがって $\Delta_Y$ の
表現可能性は $F_n$ に対する同じ主張を含意する。ゆえに、スキーム $W$ で $S$ 上のものと、
全射エタール射 $W\to F_n$ が存在することを証明すれば十分である。

\medskip\noindent
$W\to F_n$ を構成するため、エタール被覆 $\{Y_i\to Y\}$ で、$Y_i$ がスキームと
なるものを選ぶ。$X_i=X\times_YY_i$ とおき、$\mathcal{F}_i$ を
$\mathcal{F}$ の $X_i$ への引き戻しとする。すると $\mathcal{F}_i$ は、定義
または Lemma \ref{spaces-flat-lemma-quasi-finite-base-change} により $Y_i$ に関して純粋で
ある。定理の他の仮定も保たれる。最後に、$F_n$ の $Y_i$ への制限は、関手
$F_n$、すなわち $X_i\to Y_i$ と $\mathcal{F}_i$ に対応するもの、である。したがって、
次を示せば十分である。定理の主張にあるような $\mathcal{F}$ と
$f:X\to Y$ が与えられ、$Y$ がスキームであるとき、関手 $F_n$ は
スキーム $Z_n$ により表現され、$Z_n\to Y$ は有限表示の単射である。

\medskip\noindent
有限表示の単射は分離かつ準有限であることに注意せよ（Morphisms, Lemma
\ref{morphisms-lemma-monomorphism-loc-finite-type-loc-quasi-finite}）。したがって、
Descent, Lemma \ref{descent-lemma-descent-data-sheaves}、More on Morphisms,
Lemma
\ref{more-morphisms-lemma-separated-locally-quasi-finite-morphisms-fppf-descend}、
および Descent の補題 \ref{descent-lemma-descending-property-monomorphism} と
\ref{descent-lemma-descending-property-finite-presentation} を組み合わせると、
この問題は $Y$ のエタール位相に関して局所的であることが分かる。

\medskip\noindent
特に、この状況は $Y$ のザリスキー位相に関して局所的なので、$Y$ がアフィンと
仮定してよい。この場合、$f$ のファイバーの次元は上に有界であるから、$F_n$ は
十分大きい $n$ に対して表現可能である。したがって $n$ に関する降下帰納法を
使える。$F_{n+1}$ が有限表示の単射 $Z_{n+1}\to Y$ によって表現されると
仮定する。基底変換 $X_{n+1}=Z_{n+1}\times_YX$ と、引き戻し
$\mathcal{F}_{n+1}$（$\mathcal{F}$ の $X_{n+1}$ へのもの）を考える。射 $Z_{n+1}\to Y$ は
有限表示の単射なので準有限であり、Lemma
\ref{spaces-flat-lemma-quasi-finite-base-change} により $\mathcal{F}_{n+1}$ は
$Z_{n+1}$ に関して純粋である。$F_n$ は $F_{n+1}$ の部分関手なので、
$F_n$ に対する結果を示すには、状況
$\mathcal{F}_{n+1}/X_{n+1}/Z_{n+1}$ に対応する関手に対して結果を示せば
十分である。このようにして、$F_n$ に対する結果を、% P09-ERR-1265
$Y_{n+1}=Y$ の場合へ帰着する。すなわち、
$\mathcal{F}$ が次元 $\geq n+1$ において $Y$ 上平坦であると仮定してよい。

\medskip\noindent
$n$ を固定し、$\mathcal{F}$ が次元 $\geq n+1$ においてアフィンスキーム $Y$ 上
平坦であると仮定する。証明を完了するには、$F_n$ が有限表示の単射
% P09-ERR-1266
$Z_n\to S$ によって表現されることを示さなければならない。この問題は $Y$ の
エタール位相に関して局所的なので、すべての $y\in Y$ に対し、エタール近傍
$(Y',y')\to(Y,y)$ が存在し、$Y'$ への基底変換後に結果が成り立つことを示せば
十分である。
したがって Lemma \ref{spaces-flat-lemma-existence-complete} により、エタール射
$h_j:W_j\to X$（$j=1,\ldots,m$）で、各 $j$ に対して完全なデヴィサージュ
$\mathcal{F}_j/W_j/Y$ が $y$ 上に存在し、$\mathcal{F}_j$ は
$\mathcal{F}$ の $W_j$ への引き戻しであり、かつ
$|X_y|\subset\bigcup h_j(W_j)$ となるものが存在すると仮定してよい。
$h_j$ はエタールなので、Lemma \ref{spaces-flat-lemma-pre-flat-dimension-n} により、
% P09-ERR-1267
層 $\mathcal{F}_j$ は依然として次元 $\geq n+1$ において $Y$ 上平坦である。
$W=\bigcup h_j(W_j)$ とおく。これは $X$ の準コンパクト開部分である。
$\mathcal{F}$ は $X_y$ に沿って純粋なので、
% P09-ERR-1268
$$
E=\{t\in|Y|:\text{Ass}_{X_t}(\mathcal{F}_t)\subset W\}.
$$
は $y$ のすべての生成化を含む。Divisors on Spaces, Lemma
\ref{spaces-divisors-lemma-relative-assassin-constructible} により、$E$ は $Y$ の
構成可能部分集合である。上で $\Spec(\mathcal{O}_{Y,y})\subset E$ を示した。
Morphisms, Lemma \ref{morphisms-lemma-constructible-containing-open} により、
$E$ は $y$ の開近傍を含む。したがって $Y$ を縮小すれば $E=Y$ と仮定できる。
Lemma \ref{spaces-flat-lemma-localize-flat-dimension-n} により、関手 $F_n$、すなわち
$X=\coprod W_j$ および $\mathcal{F}=\coprod\mathcal{F}_j$ に対応するもの、に対して
補題を証明すれば十分である。
$F_{j,n}$ を $W_j\to Y$ と層 $\mathcal{F}_j$ に対する関手とすると、
$F_n=\prod F_{j,n}$ である。したがって、各 $F_{j,n}$ が有限表示の単射
$Z_{j,n}\to Y$ によって表現されることを示せば十分である。この場合、
$$
Z_n=Z_{1,n}\times_Y\ldots\times_YZ_{m,n}
$$
だからである。これで定理は More on Flatness, Lemma
\ref{flat-lemma-flat-dimension-n-representable} で扱った特別な場合に帰着した。
\end{proof}

\noindent
以上により、ようやく求める結果を得る。

\begin{lemma}
\label{spaces-flat-lemma-when-universal-flattening}
スキーム $S$ とする。$f:X\to Y$ を $S$ 上の代数空間の射とする。
$\mathcal{F}$ を準連接 $\mathcal{O}_X$-加群とする。
\begin{enumerate}
\item $f$ が有限表示、$\mathcal{F}$ が有限表示 $\mathcal{O}_X$-加群で、
$\mathcal{F}$ が $Y$ に関して純粋ならば、普遍平坦化 $Y'\to Y$ で
$\mathcal{F}$ のものが存在する。さらに $Y'\to Y$ は有限表示の単射である。
\item $f$ が有限表示で $X$ が $Y$ に関して純粋ならば、普遍平坦化
$Y'\to Y$ で $X$ のものが存在する。さらに $Y'\to Y$ は有限表示の単射である。
\item $f$ が固有かつ有限表示で、$\mathcal{F}$ が有限表示
$\mathcal{O}_X$-加群ならば、普遍平坦化 $Y'\to Y$ で $\mathcal{F}$ のものが
存在する。さらに $Y'\to Y$ は有限表示の単射である。
\item $f$ が固有かつ有限表示ならば、普遍平坦化 $Y'\to Y$ で $X$ のものが存在する。
\end{enumerate}
\end{lemma}

\begin{proof}
これらの主張は、Theorem \ref{spaces-flat-theorem-flat-dimension-n-representable} を
$F_0=F_{flat}$ に適用し、$f$ が固有ならば $\mathcal{F}$ は基底上自動的に
純粋であるという事実を用いれば直ちに従う。Lemma
\ref{spaces-flat-lemma-proper-pure} を参照せよ。
\end{proof}

\section{Grothendieck の存在定理}
\label{spaces-flat-section-existence}

\noindent
この節は More on Flatness, Section \ref{flat-section-existence} に対応し、
More on Morphisms of Spaces, Section
\ref{spaces-more-morphisms-section-existence-proper} の議論を続ける。
以下の状況で議論する。

\begin{situation}
\label{spaces-flat-situation-existence}
ここでは、全射な遷移写像をもち、その核が局所冪零である環の逆系 $(A_n)$ が
与えられている。$A=\lim A_n$ とおく。代数空間 $X$ が与えられ、これは
$A$ 上分離かつ有限表示である。$X_n=X\times_{\Spec(A)}\Spec(A_n)$ とおき、これを $X$ の
閉部分空間とみなす。さらに系 $(\mathcal{F}_n,\varphi_n)$ が与えられ、ここで
$\mathcal{F}_n$ は有限表示 $\mathcal{O}_{X_n}$-加群で、$A_n$ 上平坦、台が
$A_n$ 上固有であり、
$$
\varphi_n:
\mathcal{F}_n\otimes_{\mathcal{O}_{X_n}}\mathcal{O}_{X_{n-1}}
\longrightarrow
\mathcal{F}_{n-1}
$$
は同型である（Morphisms of Spaces, Lemma
\ref{spaces-morphisms-lemma-i-star-equivalence} の同値性を使った記法である）。
\end{situation}

\noindent
目標は、準連接層 $\mathcal{F}$ を $X$ 上に見つけ、等式
$\mathcal{F}_n=\mathcal{F}\otimes_{\mathcal{O}_X}\mathcal{O}_{X_n}$ がすべての $n$ に
対して成り立つようにできるかどうかを調べることである。

\begin{lemma}
\label{spaces-flat-lemma-compute-what-it-should-be}
Situation \ref{spaces-flat-situation-existence} において、
$$
K=R\lim_{D_\QCoh(\mathcal{O}_X)}(\mathcal{F}_n)=
DQ_X(R\lim_{D(\mathcal{O}_X)}\mathcal{F}_n)
$$
を考える。このとき $K$ は $D^b_{\QCoh}(\mathcal{O}_X)$ に属し、実際、$K$ の
非零コホモロジー層は次数 $\geq0$ にしか現れない。
\end{lemma}

\begin{proof}
Derived Categories of Spaces, Example
\ref{spaces-perfect-example-inverse-limit-quasi-coherent} の特別な場合である。
\end{proof}

\begin{lemma}
\label{spaces-flat-lemma-compute-against-perfect}
Situation \ref{spaces-flat-situation-existence} において、$K$ を Lemma
\ref{spaces-flat-lemma-compute-what-it-should-be} のものとする。任意の完全対象 $E$ で
$D(\mathcal{O}_X)$ に属するものに対して、次が成り立つ。
\begin{enumerate}
\item $M=R\Gamma(X,K\otimes^\mathbf{L}E)$ は $D(A)$ の完全対象であり、
標準同型
$R\Gamma(X_n,\mathcal{F}_n\otimes^\mathbf{L}E|_{X_n})=
M\otimes_A^\mathbf{L}A_n$
が $D(A_n)$ において存在する。
\item $N=R\Hom_X(E,K)$ は $D(A)$ の完全対象であり、標準同型
$R\Hom_{X_n}(E|_{X_n},\mathcal{F}_n)=N\otimes_A^\mathbf{L}A_n$
が $D(A_n)$ において存在する。
\end{enumerate}
両方の主張において、$E|_{X_n}$ は $E$ の $X_n$ への導来引き戻しを表す。
\end{lemma}

\begin{proof}
(2) の証明。$E_n=E|_{X_n}$ および
$N_n=R\Hom_{X_n}(E_n,\mathcal{F}_n)$ とおく。
$R\Hom_{X_n}(-,-)$ は $R\Gamma(X_n,R\SheafHom(-,-))$ に等しいことを
思い出そう。Cohomology on Sites, Section
\ref{sites-cohomology-section-global-RHom} を参照せよ。したがって Derived
Categories of Spaces, Lemma
\ref{spaces-perfect-lemma-base-change-RHom-perfect} により、$N_n$ は $D(A_n)$ の
完全対象であり、その形成は基底変換と可換する。ゆえに写像
$N_n\otimes_{A_n}^\mathbf{L}A_{n-1}\to N_{n-1}$ は $\varphi_n$ から得られ、
同型である。More on Algebra, Lemma
\ref{more-algebra-lemma-Rlim-perfect-gives-perfect} により、$R\lim N_n$ は完全で、
$A_n$ への基底変換により $N_n$ に戻ることが分かる。一方、三角圏の完全関手
$R\Hom_X(E,-):D_\QCoh(\mathcal{O}_X)\to D(A)$ は積と可換し、したがって
導来極限とも可換する。ゆえに
$$
R\Hom_X(E,K)=
R\lim R\Hom_X(E,\mathcal{F}_n)=
R\lim R\Hom_X(E_n,\mathcal{F}_n)=
R\lim N_n
$$
である。これで (2) が証明された。(1) は Cohomology on Sites, Lemma
\ref{sites-cohomology-lemma-dual-perfect-complex} を使って (2) に翻訳すればよい。
\end{proof}

\begin{lemma}
\label{spaces-flat-lemma-relative-pseudo-coherence}
Situation \ref{spaces-flat-situation-existence} において、$K$ を Lemma
\ref{spaces-flat-lemma-compute-what-it-should-be} のものとする。このとき $K$ は $A$ に関して
擬連接である。
\end{lemma}

\begin{proof}
% P09-ERR-1269
Lemma \ref{spaces-flat-lemma-compute-against-perfect} と Derived Categories of Spaces,
Lemma \ref{spaces-perfect-lemma-perfect-enough} を組み合わせると、
$R\Gamma(X,K\otimes^\mathbf{L}E)$ は $D(A)$ において擬連接であることが分かる。
これは、すべての擬連接対象 $E$ について成り立ち、その対象は
$D(\mathcal{O}_X)$ に属する。
したがって補題は More on Morphisms of Spaces, Lemma
\ref{spaces-more-morphisms-lemma-characterize-pseudo-coherent} から従う。
\end{proof}

\begin{lemma}
\label{spaces-flat-lemma-compute-over-affine}
Situation \ref{spaces-flat-situation-existence} において、$K$ を Lemma
\ref{spaces-flat-lemma-compute-what-it-should-be} のものとする。任意のエタール射 $U\to X$ で、
$U$ が準コンパクトかつ準分離であるものに対して、
$$
R\Gamma(U,K)\otimes_A^\mathbf{L}A_n=
R\Gamma(U_n,\mathcal{F}_n)
$$
が $D(A_n)$ において成り立つ。ここで $U_n=U\times_X X_n$ である。
\end{lemma}

\begin{proof}
$n$ を固定する。Derived Categories of Spaces, Lemma
\ref{spaces-perfect-lemma-computing-sections-as-colim} により、完全複体の系 $E_m$ で、
各項が $X$ 上にあり、
$R\Gamma(U,K)=\text{hocolim} R\Gamma(X,K\otimes^\mathbf{L}E_m)$
を満たすものが存在する。実際、この公式は $K$ に限らず、
$D_\QCoh(\mathcal{O}_X)$ のすべての対象に対して成り立つ。これを
$\mathcal{F}_n$ に適用すると、
\begin{align*}
R\Gamma(U_n,\mathcal{F}_n)
&=
R\Gamma(U,\mathcal{F}_n)\\
&=
\text{hocolim}_m R\Gamma(X,\mathcal{F}_n\otimes^\mathbf{L}E_m)\\
&=
\text{hocolim}_m R\Gamma(X_n,\mathcal{F}_n\otimes^\mathbf{L}E_m|_{X_n})
\end{align*}
を得る。Lemma \ref{spaces-flat-lemma-compute-against-perfect} と、
$-\otimes_A^\mathbf{L}A_n$ がホモトピー余極限と可換するという事実を使えば、
結論が得られる。
\end{proof}

\begin{lemma}
\label{spaces-flat-lemma-finitely-presented}
Situation \ref{spaces-flat-situation-existence} において、$K$ を Lemma
\ref{spaces-flat-lemma-compute-what-it-should-be} のものとする。$X_0\subset|X|$ を、閉部分集合
$\Spec(A_1)=\Spec(A_2)=\ldots$ の上にある点からなる閉部分集合と記す。この等式の
各項は $\Spec(A)$ の閉部分集合とみなしている。
% P09-ERR-1271
開部分空間 $W\subset X$ で $X_0$ を含み、次を満たすものが存在する。
\begin{enumerate}
\item $H^i(K)|_W$ は $i=0$ でない限り零である。
\item $\mathcal{F}=H^0(K)|_W$ は有限表示である。
\item $\mathcal{F}_n=\mathcal{F}\otimes_{\mathcal{O}_X}\mathcal{O}_{X_n}$ である。
\end{enumerate}
\end{lemma}

\begin{proof}
$n\geq1$ を固定する。構成により標準写像 $K\to\mathcal{F}_n$ が
$D_\QCoh(\mathcal{O}_X)$ にあり、したがって準連接層の標準写像
$H^0(K)\to\mathcal{F}_n$ がある。これにより (3) の意味が明らかになる。

\medskip\noindent
点 $x\in X_0$ を取る。(1)、(2)、(3) が成り立つ開近傍 $W$ で $x$ を含むものを見つける。
$X_0$ は準コンパクトなので、これで補題が証明される。エタール射 $U\to X$ を取り、
$U$ はアフィンで、点 $u\in U$ は $x$ に写るとする。写像 $|U|\to|X|$ は開なので、
$u$ の $U$ 内の開近傍で (1)、(2)、(3) が成り立つものを見つければ十分である。
$U=\Spec(B)$ と書く。$P\to B$ なる全射で、$P$ が $A$ 上滑らかなものを選ぶ。
Lemma \ref{spaces-flat-lemma-relative-pseudo-coherence} と相対擬連接性の定義により、上に有界な
複体 $F^\bullet$ で、有限自由 $P$-加群からなり、$Ri_*K$ を表すものが存在する。
ここで $i:U\to\Spec(P)$ は表示から誘導される閉埋め込みである。$M_n$ を $B$-加群で、
$\mathcal{F}_n|_U$ に対応するものとする。Lemma
\ref{spaces-flat-lemma-compute-over-affine} により、
$$
H^i(F^\bullet\otimes_A A_n)=
\left\{
\begin{matrix}
0&\text{if}&i\not=0\\
M_n&\text{if}&i=0
\end{matrix}
\right.
$$
添字 $i$ を、$F^i$ が非零となるもののうち最大のものとする。$i\leq0$ ならば (1)、(2)、(3) は
成り立つ。そうでなければ $i>0$ であり、写像
$$
F^{i-1}\to F^i
$$
の点 $u$ における階数が最大であることが分かる。したがって $u$ の $\Spec(P)$ における開近傍でも
階数は最大である。そこで $P$ を主局所化で置き換え、表示された写像が全射であると
仮定してよい。$F^i$ は有限自由なので、分裂 $F^{i-1}=F'\oplus F^i$ を選べる。
このとき $F^\bullet$ を複体
$$
\ldots\to F^{i-2}\to F'\to0\to\ldots
$$
で置き換えられ、$i$ に関する帰納法で結論を得る。
\end{proof}

\begin{lemma}
\label{spaces-flat-lemma-proper-support}
Situation \ref{spaces-flat-situation-existence} において、$K$ を Lemma
\ref{spaces-flat-lemma-compute-what-it-should-be} のものとする。$W\subset X$ を Lemma
\ref{spaces-flat-lemma-finitely-presented} のものとし、$\mathcal{F}=H^0(K)|_W$ とおく。
このとき、必要なら開集合 $W$ を縮小することにより、$\mathcal{F}$ の台は $A$ 上
固有となる。
\end{lemma}

\begin{proof}
$n\geq1$ を固定する。$I_n=\Ker(A\to A_n)$ とおく。More on Algebra, Lemma
\ref{more-algebra-lemma-limit-henselian} により、対 $(A,I_n)$ は Hensel 対である。
$Z\subset W$ を $\mathcal{F}$ のスキーム論的台とする。$\mathcal{F}$ は有限表示なので、
これは閉部分空間である。Lemma \ref{spaces-flat-lemma-finitely-presented} の (3) により、
$Z\times_{\Spec(A)}\Spec(A_n)$ は $\mathcal{F}_n$ の台に等しく、したがって
% P09-ERR-1272
$\Spec(A/I)$ 上固有である。More on Morphisms of Spaces, Lemma
\ref{spaces-more-morphisms-lemma-split-off-proper-part-henselian} により、
$Z=Z_1\amalg Z_2$ と書ける。ここで $Z_1,Z_2$ は $Z$ の開かつ閉な部分空間で、
$Z_1$ は $A$ 上固有であり、$Z_1\times_{\Spec(A)}\Spec(A/I_n)$ は
$\mathcal{F}_n$ の台に等しい。言い換えると、$|Z_2|$ は $X_0$ と交わらない。
したがって $W$ を $W\setminus Z_2$ で置き換えれば、補題が得られる。
\end{proof}

\begin{theorem}[Grothendieck の存在定理]
\label{spaces-flat-theorem-existence}
Situation \ref{spaces-flat-situation-existence} において、有限表示 $\mathcal{O}_X$-加群
$\mathcal{F}$ で、$A$ 上平坦かつ台が $A$ 上固有であり、等式
$\mathcal{F}_n=\mathcal{F}\otimes_{\mathcal{O}_X}\mathcal{O}_{X_n}$ がすべての $n$ に対して
写像 $\varphi_n$ と両立するように成り立つものが存在する。
\end{theorem}

\begin{proof}
Lemmas \ref{spaces-flat-lemma-compute-what-it-should-be},
\ref{spaces-flat-lemma-compute-against-perfect},
\ref{spaces-flat-lemma-relative-pseudo-coherence},
\ref{spaces-flat-lemma-compute-over-affine},
\ref{spaces-flat-lemma-finitely-presented}, および
\ref{spaces-flat-lemma-proper-support} を適用すると、開部分空間 $W\subset X$ で、
$\Spec(A_n)$ の上にあるすべての点を含むものと、有限表示 $\mathcal{O}_W$-加群
$\mathcal{F}$ で、台が $A$ 上固有かつ、等式
$\mathcal{F}_n=\mathcal{F}\otimes_{\mathcal{O}_W}\mathcal{O}_{X_n}$ をすべての
$n\geq1$ に対して満たすものが得られる
（$X_n\subset W$ なので、これは意味をもつ）。Lemma \ref{spaces-flat-lemma-proper-pure} により、
$\mathcal{F}$ は $\Spec(A)$ に関して普遍的に純粋である。Theorem
\ref{spaces-flat-theorem-flat-dimension-n-representable}（説明については Lemma
\ref{spaces-flat-lemma-when-universal-flattening} を参照）により、普遍平坦化
$S'\to\Spec(A)$ が存在する。これは $\mathcal{F}$ に対するものであり、さらに射
$S'\to\Spec(A)$ は有限表示の単射である。
特に $S'$ はスキームである（これは定理の証明から従うが、
% P09-ERR-1273
事後的には Morphisms of Spaces, Proposition
\ref{spaces-morphisms-proposition-locally-quasi-finite-separated-over-scheme}
からも従う）。$\mathcal{F}$ の $\Spec(A_n)$ への基底変換は $\mathcal{F}_n$ なので、
射 $\Spec(A_n)\to\Spec(A)$ は $S'$ を経由して（一意に）分解し、これは各 $n$ に
対して成り立つ。
More on Flatness, Lemma \ref{flat-lemma-monomorphism-isomorphism} により、
$S'=\Spec(A)$ であることが分かる。これは $\mathcal{F}$ が $A$ 上平坦であることを
意味する。最後に、スキーム論的台 $Z$ は $\mathcal{F}$ の台であり、$\Spec(A)$ 上
固有なので、
射 $Z\to X$ は閉である。したがって、
% P09-ERR-1274
押し出し $(W\to X)_*\mathcal{F}$ は $W$ 上に台をもち、所望の性質をすべて満たす。
\end{proof}







\section{Grothendieck の存在定理、再論}
\label{spaces-flat-section-existence-derived}

\noindent
この節では、準連接加群について Section \ref{spaces-flat-section-existence} で用いた方法に従い、
導来圏における Grothendieck の存在定理の類似を証明する。この節は代数空間に対する
More on Flatness, Section \ref{flat-section-existence-derived} の類似である。古典的な場合
（代数空間に対するもの）は More on Morphisms of Spaces, Section
\ref{spaces-more-morphisms-section-existence-proper} で論じられている。以下の状況で議論する。

\begin{situation}
\label{spaces-flat-situation-existence-derived}
ここでは、全射な遷移写像をもち、その核が局所冪零である環の逆系 $(A_n)$ が
与えられている。$A=\lim A_n$ とおく。代数空間 $X$ があり、これは $A$ 上固有、
平坦かつ有限表示である。$X_n=X\times_{\Spec(A)}\Spec(A_n)$ とおき、これを $X$ の
閉部分空間とみなす。さらに系 $(K_n,\varphi_n)$ が与えられ、ここで $K_n$ は
$D(\mathcal{O}_{X_n})$ の擬連接対象であり、
$$
\varphi_n:K_n\longrightarrow K_{n-1}
$$
は $D(\mathcal{O}_{X_n})$ における写像で、同型
$K_n\otimes_{\mathcal{O}_{X_n}}^\mathbf{L}\mathcal{O}_{X_{n-1}}
\to K_{n-1}$ を $D(\mathcal{O}_{X_{n-1}})$ において誘導する。
\end{situation}

\noindent
より正確には、
$\varphi_n:K_n\to Ri_{n-1,*}K_{n-1}$
と書くべきである。ここで $i_{n-1}:X_{n-1}\to X_n$ は包含射であり、この記法では、
随伴写像 $Li_{n-1}^*K_n\to K_{n-1}$ が同型であることが条件となる。目標は、擬連接な
$K\in D(\mathcal{O}_X)$ で、等式
$K_n=K\otimes_{\mathcal{O}_X}^\mathbf{L}\mathcal{O}_{X_n}$ をすべての $n$ に対して
満たすものを見つけることである
（同じ記法上の濫用を用いる）。

\begin{lemma}
\label{spaces-flat-lemma-compute-what-it-should-be-derived}
Situation \ref{spaces-flat-situation-existence-derived} において、
$$
K=R\lim_{D_\QCoh(\mathcal{O}_X)}(K_n)=
DQ_X(R\lim_{D(\mathcal{O}_X)}K_n)
$$
を考える。このとき $K$ は $D^-_{\QCoh}(\mathcal{O}_X)$ に属する。
\end{lemma}

\begin{proof}
関手 $DQ_X$ は、$X$ が準コンパクトかつ準分離なので存在する。Derived Categories of
Spaces, Lemma \ref{spaces-perfect-lemma-better-coherator} を参照せよ。$DQ_X$ は右随伴なので
積と可換し、したがって導来極限とも可換する。これにより補題の主張にある等式を得る。

\medskip\noindent
Derived Categories of Spaces, Lemma
\ref{spaces-perfect-lemma-boundedness-better-coherator} により、関手 $DQ_X$ は有界な
コホモロジー次元をもつ。したがって $R\lim K_n\in D^-(\mathcal{O}_X)$ を示せば十分である。
これを見るため、エタール射 $U\to X$ で $U$ がアフィンなものを取る。このとき標準完全列
$$
0\to
R^1\lim H^{m-1}(U,K_n)\to H^m(U,R\lim K_n)\to
\lim H^m(U,K_n)\to0
$$
が Cohomology on Sites, Lemma
\ref{sites-cohomology-lemma-RGamma-commutes-with-Rlim} により存在する。$U$ はアフィンで、
$K_n$ は擬連接である（したがって Derived Categories of Spaces, Lemma
\ref{spaces-perfect-lemma-pseudo-coherent} により準連接コホモロジー層をもつ）ので、
$H^m(U,K_n)=H^m(K_n)(U)$ が Derived Categories of Schemes, Lemma
\ref{perfect-lemma-affine-compare-bounded} により成り立つ。ゆえに $K_n$ が $n$ に依存せず
一様に上に有界であることを示せば十分である。

\medskip\noindent
$K_n$ は擬連接なので、$K_n\in D^-(\mathcal{O}_{X_n})$ である。$a_n$ を
$H^{a_n}(K_n)$ が非零となる最大の整数とする。もちろん
$a_1\leq a_2\leq a_3\leq\ldots$ である。$H^{a_n}(K_n)$ は有限表示
$\mathcal{O}_{X_n}$-加群であることに注意せよ（Cohomology on Sites, Lemma
\ref{sites-cohomology-lemma-finite-cohomology}）。等式
$H^{a_n}(K_{n-1})=
H^{a_n}(K_n)\otimes_{\mathcal{O}_{X_n}}\mathcal{O}_{X_{n-1}}$
が成り立つ。$X_{n-1}\to X_n$ は厚化なので、Nakayama の補題（Algebra, Lemma
\ref{algebra-lemma-NAK}）により、
$H^{a_n}(K_n)\otimes_{\mathcal{O}_{X_n}}\mathcal{O}_{X_{n-1}}$
が零ならば $H^{a_n}(K_n)$ も零である（例えば茎で確認すればよい。細部は省略する）。
したがって $a_{n-1}=a_n$ であり、これはすべての $n$ に対して成り立つ。結論を得た。
\end{proof}

\begin{lemma}
\label{spaces-flat-lemma-compute-against-perfect-derived}
Situation \ref{spaces-flat-situation-existence-derived} において、$K$ を Lemma
\ref{spaces-flat-lemma-compute-what-it-should-be-derived} のものとする。任意の完全対象 $E$ で
$D(\mathcal{O}_X)$ に属するものについて、コホモロジー
$$
M=R\Gamma(X,K\otimes^\mathbf{L}E)
$$
は $D(A)$ の擬連接対象であり、標準同型
$$
R\Gamma(X_n,K_n\otimes^\mathbf{L}E|_{X_n})=M\otimes_A^\mathbf{L}A_n
$$
が $D(A_n)$ において存在する。ここで $E|_{X_n}$ は $E$ の $X_n$ への導来引き戻しを表す。
\end{lemma}

\begin{proof}
$E_n=E|_{X_n}$ および
$M_n=R\Gamma(X_n,K_n\otimes^\mathbf{L}E|_{X_n})$ とおく。Derived Categories of
Spaces, Lemma
\ref{spaces-perfect-lemma-flat-proper-pseudo-coherent-direct-image-general} により、$M_n$ は
$D(A_n)$ の擬連接対象で、その形成は基底変換と可換する。したがって写像
$M_n\otimes_{A_n}^\mathbf{L}A_{n-1}\to M_{n-1}$ は $\varphi_n$ から得られ、同型である。
More on Algebra, Lemma
\ref{more-algebra-lemma-Rlim-pseudo-coherent-gives-pseudo-coherent} により、$R\lim M_n$ は
擬連接であり、$A_n$ への基底変換により $M_n$ に戻る。一方、三角圏の完全関手
$R\Gamma(X,-):D_\QCoh(\mathcal{O}_X)\to D(A)$ は積と可換し、したがって導来極限とも
可換する。ゆえに
$$
R\Gamma(X,E\otimes^\mathbf{L}K)=
R\lim R\Gamma(X,E\otimes^\mathbf{L}K_n)=
R\lim R\Gamma(X_n,E_n\otimes^\mathbf{L}K_n)=
R\lim M_n
$$
となり、所望の結論を得る。
\end{proof}

\begin{lemma}
\label{spaces-flat-lemma-relative-pseudo-coherence-derived}
Situation \ref{spaces-flat-situation-existence-derived} において、$K$ を Lemma
\ref{spaces-flat-lemma-compute-what-it-should-be-derived} のものとする。このとき $K$ は $X$ 上
擬連接である。
\end{lemma}

\begin{proof}
% P09-ERR-1275
Lemma \ref{spaces-flat-lemma-compute-against-perfect-derived} と Derived Categories of Spaces, Lemma
\ref{spaces-perfect-lemma-perfect-enough} を組み合わせると、
$R\Gamma(X,K\otimes^\mathbf{L}E)$ は $D(A)$ において擬連接であることが分かる。
これは、すべての擬連接対象 $E$ について成り立ち、その対象は $D(\mathcal{O}_X)$ に
属する。したがって More on Morphisms of Spaces, Lemma
\ref{spaces-more-morphisms-lemma-characterize-pseudo-coherent} により、$K$ は $A$ に関して
擬連接である。
% P09-ERR-1276
$X$ は $A$ 上平坦かつ有限表示なので、これは $X$ 上擬連接であることと同値である。
More on Morphisms of Spaces, Lemma
\ref{spaces-more-morphisms-lemma-relative-pseudo-coherent-is-moot} を参照せよ。
\end{proof}

\begin{lemma}
\label{spaces-flat-lemma-compute-over-affine-derived}
Situation \ref{spaces-flat-situation-existence-derived} において、$K$ を Lemma
\ref{spaces-flat-lemma-compute-what-it-should-be-derived} のものとする。任意のエタール射 $U\to X$ で、
$U$ が準コンパクトかつ準分離であるものに対して、
$$
R\Gamma(U,K)\otimes_A^\mathbf{L}A_n=
R\Gamma(U_n,K_n)
$$
が $D(A_n)$ において成り立つ。ここで $U_n=U\times_X X_n$ である。
\end{lemma}

\begin{proof}
$n$ を固定する。Derived Categories of Spaces, Lemma
\ref{spaces-perfect-lemma-computing-sections-as-colim} により、完全複体の系 $E_m$ で、各項が
$X$ 上にあり、
$R\Gamma(U,K)=\text{hocolim} R\Gamma(X,K\otimes^\mathbf{L}E_m)$
を満たすものが存在する。実際、この公式は $K$ に限らず、
$D_\QCoh(\mathcal{O}_X)$ のすべての対象に対して成り立つ。これを $K_n$ に適用すると、
\begin{align*}
R\Gamma(U_n,K_n)
&=
R\Gamma(U,K_n)\\
&=
\text{hocolim}_m R\Gamma(X,K_n\otimes^\mathbf{L}E_m)\\
&=
\text{hocolim}_m R\Gamma(X_n,K_n\otimes^\mathbf{L}E_m|_{X_n})
\end{align*}
を得る。Lemma \ref{spaces-flat-lemma-compute-against-perfect-derived} と、
$-\otimes_A^\mathbf{L}A_n$ がホモトピー余極限と可換するという事実を使えば、結論が得られる。
\end{proof}

\begin{theorem}[導来 Grothendieck 存在定理]
\label{spaces-flat-theorem-existence-derived}
Situation \ref{spaces-flat-situation-existence-derived} において、擬連接な $K$ で
$D(\mathcal{O}_X)$ に属し、等式
$K_n=K\otimes_{\mathcal{O}_X}^\mathbf{L}\mathcal{O}_{X_n}$ がすべての $n$ に対して
写像 $\varphi_n$ と両立するように成り立つものが存在する。
\end{theorem}

\begin{proof}
Lemmas \ref{spaces-flat-lemma-compute-what-it-should-be-derived},
\ref{spaces-flat-lemma-compute-against-perfect-derived},
\ref{spaces-flat-lemma-relative-pseudo-coherence-derived} を適用すると、擬連接対象 $K$ で
$D(\mathcal{O}_X)$ に属するものが得られる。Lemma
\ref{spaces-flat-lemma-compute-over-affine-derived} でアフィンな $U$ を選ぶと、$K$ の制限は
$K_n$ に一致し、その制限先は $X_n$ であることが直ちに従う。
\end{proof}

\begin{remark}
\label{spaces-flat-remark-correct-generality}
この節の結果は一般化できる。$X\to\Spec(A)$ が分離かつ有限表示であり、$K_n$ が
$A_n$ に関して擬連接で、$X_n$ の閉部分集合上に台をもち、その閉部分集合が $A_n$ 上
固有であることだけを仮定しても、おそらく正しい。得られる $K$ は $A$ に関して擬連接で、
$A$ 上固有な閉部分集合に台をもつ。この結果が必要になった時点で、正確な主張を定式化し、
ここで証明することにする。
\end{remark}
